% ============================================================================
% Część 2 przewodnika: podręcznik obsługi (prefiks diagramów: p2).
% Wczytywana przez \input z przewodnik.pl.tex -- zawiera wyłącznie treść od \section.
%
% Konwencja dowodowa tej części: każde wyjście w ramkach Verbatim zostało
% zebrane na żywo w sesji pisania (2026-10-05, Linux, .NET SDK 10.0.112,
% Node 22.22.0, pnpm 12.9.1) i skrócone tylko tam, gdzie widać „...”.
% Czego w tej sesji nie dało się uruchomić (Docker/Aspire, obraz authservice,
% prawdziwy model, Claude), oznaczono wprost „nieprzetestowane w tej sesji”.
% ============================================================================

{\sloppy\emergencystretch=4em
\section{Wymagania i instalacja}

Ta część jest podręcznikiem obsługi: pokazuje, co wpisać, czego się spodziewać
i co zrobić, gdy wynik jest inny. Wszystkie polecenia uruchamia się z katalogu
głównego repozytorium (\repopath{exit-interview-agent}). Każde wyjście w ramce
pochodzi z faktycznego uruchomienia; jeśli czegoś nie dało się uruchomić w
środowisku, w którym powstał ten tekst, napisano to wprost jako
\emph{nieprzetestowane w tej sesji} i podano plik z dokumentacją, na której
opis się opiera.

\subsection{Co musi być zainstalowane}

\begin{center}
\small
\begin{tabularx}{\linewidth}{@{}l>{\raggedright\arraybackslash}X>{\raggedright\arraybackslash}X@{}}
\toprule
Narzędzie & Wymaganie (źródło) & Sprawdzone w tej sesji \\
\midrule
git & dowolny; skrypt \repopath{scripts/setup.sh} tylko sprawdza obecność & jest \\
.NET SDK & 10.x; \repopath{global.json}: \texttt{10.0.100}, \texttt{rollForward:\allowbreak{}latestFeature} & \texttt{10.0.112} \\
Node.js & 22 lub nowszy (\repopath{web/package.json}: \texttt{engines.\allowbreak{}node >=22}) & \texttt{v22.22.0} \\
pnpm & wersja z \texttt{packageManager} w \repopath{web/package.json}: \texttt{pnpm@12.9.1}; instalacja: \texttt{npm install -g pnpm} albo \texttt{corepack enable} & \texttt{12.9.1} \\
Silnik kontenerów & Docker albo Podman; potrzebny AppHost-owi (Postgres i obraz authservice) oraz budowie obrazów & \textbf{brak działającego demona} --- patrz niżej \\
\bottomrule
\end{tabularx}
\end{center}

Silnik kontenerów jest jedynym wymaganiem dla pełnego stosu Aspire. Samo CLI,
testy .NET, harness ewaluacyjny i portal webowy (z atrapą backendu w testach
e2e) działają bez niego --- w tej sesji właśnie tak je uruchomiono.

\subsection{Skrypt przygotowawczy}

Na Linuksie i macOS: \komenda{scripts/\allowbreak{}setup.\allowbreak{}sh}, na Windows:
\komenda{scripts/\allowbreak{}setup.\allowbreak{}ps1}. Oba mają tryb tylko-raport (\komenda{--check} /
\komenda{-Check}), który niczego nie zmienia. Skrypt wykonuje cztery numerowane
kroki (z komentarzy w \repopath{scripts/setup.sh}):

\begin{enumerate}[leftmargin=*,itemsep=2pt]
  \item sprawdza wymagania wstępne (tabela wyżej; \texttt{dotnet} musi mieć
        wersję \texttt{10.*});
  \item ustawia \texttt{core.\allowbreak{}hooksPath = scripts/\allowbreak{}hooks}, czyli skanowanie
        sekretów (gitleaks) przed każdym commitem;
  \item zakłada lokalny magazyn sekretów \texttt{dotnet user-\allowbreak{}secrets} dla
        projektu AppHost i, jeśli jest \texttt{openssl}, generuje trzy
        sekrety \emph{wyłącznie deweloperskie}: klucz RSA-2048 do podpisu JWT
        (\texttt{Parameters:\allowbreak{}authservice-\allowbreak{}jwt-\allowbreak{}private-\allowbreak{}key}), sekret klienta MCP
        (\texttt{Parameters:\allowbreak{}authservice-\allowbreak{}mcp-\allowbreak{}client-\allowbreak{}secret}, 32 losowe bajty
        hex) i klucz szyfrowania tokenów authservice
        (\texttt{Parameters:\allowbreak{}authservice-\allowbreak{}encryption-\allowbreak{}key}, base64 z 32 bajtów);
  \item wypisuje, że integracje opcjonalne (prawdziwy model) nie są tu
        konfigurowane, i podaje komendę startu stosu.
\end{enumerate}

Przebieg trybu raportu, uruchomiony na maszynie bez działającego Dockera:

\begin{Verbatim}[fontsize=\scriptsize,frame=single,framesep=2mm]
$ scripts/setup.sh --check
1/4 prerequisites
  ok       git
  ok       dotnet
  ok       node
  ok       pnpm
  MISSING  container engine -> https://docs.docker.com/get-docker/ (needed by the AppHost ...)

check: 1 item(s) missing
\end{Verbatim}

Kod wyjścia to 1, gdy czegoś brakuje, i 0, gdy wszystko jest. Pełny przebieg
(bez \komenda{--check}) kończy się wskazówką
\komenda{dotnet run --project src/\allowbreak{}ExitInterviewAgent.\allowbreak{}AppHost} i odsyła do
\repopath{scripts/README.md} po tabelę rozwiązywania problemów. Pełny przebieg
\emph{nie był uruchamiany} w tej sesji, bo wymaga silnika kontenerów już w
kroku pierwszym (nieprzetestowane w tej sesji; opis wg
\repopath{scripts/setup.sh}). Sekret nigdy nie trafia do plików w drzewie
repozytorium --- tak stanowi komentarz w skrypcie i
\repopath{secrets.env.example}.

\subsection{Budowa i testy}

\begin{Verbatim}[fontsize=\scriptsize,frame=single,framesep=2mm]
$ dotnet build -warnaserror
Build succeeded.
    0 Warning(s)
    0 Error(s)

$ dotnet test
Passed!  - Failed: 0, Passed: 110 ... Records.Tests.dll      (podobnie Personas 18, Agent 401,
Privacy 184, Providers 155 [2 pominięte], Cli 130, Eval 161, Signals 102,
InterviewService 359 [15 pominiętych])

$ scripts/check-kernel-size.sh
ServiceDefaults: 320 lines (limit 800)
\end{Verbatim}

Razem 1620 testów zaliczonych i 17 pominiętych. Pominięte to testy, które mają
własny warunek uruchomienia i w tej sesji go nie spełniły: testy na prawdziwym
PostgreSQL (\texttt{PostgresSubmissionTests}, \texttt{PostgresSignals}) czekają
na zmienną \texttt{TEST\_\allowbreak{}POSTGRES\_\allowbreak{}CONNECTION}, a dwa testy
\texttt{LiveSmoke} dostawców czekają na klucz i
\texttt{EXIT\_\allowbreak{}INTERVIEW\_\allowbreak{}LIVE\_\allowbreak{}MODEL} (README, „Run it with your own model”).
Pominięcie nie jest zaliczeniem: to jest to, czego ta sesja nie sprawdziła.

Portal webowy (\komenda{cd web}):

\begin{Verbatim}[fontsize=\scriptsize,frame=single,framesep=2mm]
$ pnpm install --frozen-lockfile    # Done in 8.2s
$ pnpm lint && pnpm typecheck && pnpm test
 Test Files  18 passed (18)
      Tests  253 passed (253)
$ pnpm build                        # Next.js 16.3.8, trasy dynamiczne
\end{Verbatim}

Testy przeglądarkowe (Playwright, \repopath{tests/e2e}) uruchamiają
\emph{zbudowany} portal przeciwko atrapie backendu
(\repopath{tests/e2e/support/stub-backend.mjs}). Gdy w systemie jest już
Chromium, wskazuje się go zmienną \texttt{PLAYWRIGHT\_\allowbreak{}CHROMIUM\_\allowbreak{}EXECUTABLE}
zamiast \komenda{npx playwright install}:

\begin{Verbatim}[fontsize=\scriptsize,frame=single,framesep=2mm]
$ cd tests/e2e && pnpm install --frozen-lockfile
$ PLAYWRIGHT_CHROMIUM_EXECUTABLE=/opt/pw-browsers/chromium-1194/chrome-linux/chrome pnpm test
  115 passed (1.1m)
\end{Verbatim}

\uwaga{Testy e2e sprawdzają portal względem \emph{atrapy} serwisu, która
odzwierciedla kontrakt, a nie względem prawdziwego interview-service i
authservice (\repopath{tests/e2e/README.md}). Zielone e2e nie dowodzi, że
pełny stos Aspire działa --- tego w tej sesji nie sprawdzono.}

\section{Uruchomienie lokalne}

\subsection{Jedno polecenie: AppHost}

\begin{Verbatim}[fontsize=\scriptsize,frame=single,framesep=2mm]
dotnet run --project src/ExitInterviewAgent.AppHost
\end{Verbatim}

AppHost (\texttt{src/\allowbreak{}ExitInterviewAgent.\allowbreak{}AppHost/\allowbreak{}Program.\allowbreak{}cs}) to korzeń
kompozycji tylko do developmentu; produkcję opisują pliki
\repopath{flyio/*.fly.toml}. Z kodu wynika następująca topologia (diagram
poniżej):

\begin{center}
\includegraphics[width=0.92\linewidth,height=0.4\textheight,keepaspectratio]{\dgm{p2-stos-lokalny}}
\end{center}

\begin{itemize}[leftmargin=*,itemsep=1pt]
  \item \texttt{postgres}: jeden serwer \texttt{postgres:17-alpine}, wolumen \texttt{exit-interview-agent-pgdata}; wersja zgodna z \repopath{flyio/postgres.fly.toml}. Dwie bazy logiczne, \texttt{interviewdb} i \texttt{authdb} (w developmencie na koncie superużytkownika).
  \item \texttt{interview-service}: projekt .NET, \texttt{DATABASE\_PROVIDER=PostgreSQL}, czeka na \texttt{interviewdb}; sonda \texttt{/health}.
  \item \texttt{authservice}: \emph{osobna instancja} obrazu \texttt{ghcr.io/konradcinkusz/authservice:v0.3.4}; port hosta \textbf{5100} (stały: to adres JWKS i issuera), kontenera 8080; \texttt{Database\_\_SchemaMode=Migrate}; sonda \texttt{/health/ready}.
  \item \texttt{web}: Next.js uruchamiany jako \texttt{dev} przez pnpm (\texttt{web/app}), port ze zmiennej \texttt{PORT}, \texttt{SESSION\_COOKIE\_SECURE=false} (tylko http na localhost); sonda \texttt{/healthz}.
\end{itemize}

Panel Aspire nasłuchuje według \repopath{Properties/launchSettings.json} na
\texttt{http:\allowbreak{}/\allowbreak{}/\allowbreak{}localhost:\allowbreak{}15200}; AppHost wypisuje jego adres przy starcie.
Porty \texttt{web} i \texttt{interview-\allowbreak{}service} nadaje Aspire (README: „web and
interview-service get their ports from it”), dlatego adresów nie przepisuje
się z pamięci --- czyta się je z panelu. Gdy projekt \texttt{interview-\allowbreak{}service}
uruchamia się \emph{samodzielnie}, jego profil startowy używa
\texttt{http:\allowbreak{}/\allowbreak{}/\allowbreak{}localhost:\allowbreak{}5200}.

\uwaga{\textbf{Nieprzetestowane w tej sesji:} uruchomienie AppHost-a. Maszyna,
na której powstał ten tekst, nie ma działającego demona kontenerów
(\komenda{scripts/\allowbreak{}setup.\allowbreak{}sh --check} raportuje brak silnika), więc nie startował
ani Postgres, ani obraz authservice. Opis powyżej pochodzi z czytania
\texttt{Program.cs}, \repopath{README.md} i \repopath{scripts/README.md}.}

\subsection{Co się dzieje bez poświadczeń i bez obrazu}

Zasada projektu (P8): świeży klon bez żadnych poświadczeń uruchamia się z
ograniczonymi funkcjami, a nie przestaje działać. README podaje, że
\texttt{GET <interview-\allowbreak{}service>/\allowbreak{}health} wypisuje, co jest zdegradowane, a
baner startowy drukuje tę samą listę. Tak wygląda to na żywo --- samodzielnie
uruchomiony interview-service (bez bazy, bez tożsamości; w tej sesji w trybie
InMemory):

\begin{Verbatim}[fontsize=\scriptsize,frame=single,framesep=2mm]
$ ASPNETCORE_URLS=http://localhost:5200 \
    dotnet run --project src/ExitInterviewAgent.InterviewService
info: Startup[0] integration identity: DEGRADED (Jwt:Authority not set; protected endpoints answer 401)
info: Startup[0] integration mcp-auth: DEGRADED (Mcp:Issuer and Mcp:Resource not set ...)
info: Startup[0] integration database: DEGRADED (InMemory ...: data is lost on restart)
info: Startup[0] integration ledger-key / employment-verifier: DEGRADED (ephemeral key / mock verifier)
info: Startup[0] integration signals: configured (snapshots every 24 h ..., k = 5, rules v1 ...)

$ curl -s localhost:5200/health
{"status":"Degraded","checks":{"self":"Healthy","schema":"Healthy","signals":"Healthy"},
 "integrations":[{"name":"identity","configured":false,...}, ... ]}
\end{Verbatim}

Status \texttt{Degraded} przy poprawnych kontrolach \texttt{checks} jest stanem
normalnym w developmencie: znaczy „działa, ale integracja X nie jest
skonfigurowana”. Dalsze próby pokazują, że zabezpieczenia trzymają bez
konfiguracji:

\begin{Verbatim}[fontsize=\scriptsize,frame=single,framesep=2mm]
$ curl -s -i localhost:5200/api/v1/me         ->  HTTP/1.1 401 Unauthorized
$ curl -s -i localhost:5200/api/v1/signals/employers  ->  401, WWW-Authenticate: Bearer
$ curl -s -i -X POST localhost:5200/mcp       ->  HTTP/1.1 404 Not Found
\end{Verbatim}

Ostatnia linia jest zgodna z dokumentacją: bez \texttt{Mcp\_\_Issuer} i
\texttt{Mcp\_\_Resource} ścieżka MCP jest wyłączona.

Jeśli obrazu \texttt{ghcr.\allowbreak{}io/\allowbreak{}konradcinkusz/\allowbreak{}authservice} nie da się pobrać
(README: tak bywa za proxy), uruchamia się stos bez tożsamości:

\begin{Verbatim}[fontsize=\scriptsize,frame=single,framesep=2mm]
Identity__Enabled=false dotnet run --project src/ExitInterviewAgent.AppHost
\end{Verbatim}

Punkty chronione odpowiadają wtedy 401 i mówią o tym (README;
\texttt{docs/\allowbreak{}adr/\allowbreak{}0003-\allowbreak{}identity-\allowbreak{}authservice-\allowbreak{}as-\allowbreak{}pinned-\allowbreak{}image.\allowbreak{}md}); logowanie i
rejestracja w portalu nie działają
(\texttt{identity\_\allowbreak{}not\_\allowbreak{}configured}, patrz sekcja~\ref{sec:portal}).

\subsection{Konfiguracja i sekrety}

Pełna lista zmiennych z podziałem na warstwy (\texttt{always}, \texttt{mode},
\texttt{secret}, \texttt{ci}, \texttt{tuning}) jest w
\repopath{secrets.env.example}; autorytatywna wersja w
\repopath{scripts/README.md}. W trybie AppHost większość wartości wstrzykuje
AppHost; plik \texttt{.env} jest potrzebny tylko, gdy uruchamia się usługę poza
AppHost-em. Najczęściej potrzebne:

\begin{center}
\small
\begin{longtable}{@{}>{\raggedright\arraybackslash}p{0.38\linewidth}>{\raggedright\arraybackslash}p{0.58\linewidth}@{}}
\toprule
Klucz & Znaczenie i zachowanie bez niego \\
\midrule
\endhead
\texttt{DATABASE\_\allowbreak{}PROVIDER} & \texttt{PostgreSQL}; pusty lub brak łańcucha połączenia = InMemory (dane znikają po restarcie; \texttt{/health}: \texttt{database=inmemory}) \\
\texttt{Jwt\_\allowbreak{}\_\allowbreak{}Authority} & publiczny adres authservice; bez niego endpointy chronione dają 401 \\
\texttt{Ledger\_\allowbreak{}\_\allowbreak{}ActiveKeyId}, \texttt{Ledger\_\allowbreak{}\_\allowbreak{}Keys\_\allowbreak{}\_\allowbreak{}0\_\allowbreak{}\_\allowbreak{}Id}, \texttt{Ledger\_\allowbreak{}\_\allowbreak{}Keys\_\allowbreak{}\_\allowbreak{}0\_\allowbreak{}\_\allowbreak{}Secret} & klucz HMAC rejestru zgłoszeń (id 1--32 znaków \texttt{[A-Za-z0-9\_-]}, sekret base64 co najmniej 32 bajty). \textbf{Poza Development usługa odmawia startu bez niego}; w Development generuje się klucz tymczasowy na uruchomienie \\
\texttt{Submission\_\allowbreak{}\_\allowbreak{}ClientIpHeader} & nagłówek z prawdziwym adresem klienta, \emph{tylko} za zaufanym proxy (np. \texttt{Fly-Client-IP}); bez niego używany jest adres gniazda \\
\texttt{INTERVIEW\_\allowbreak{}SERVICE\_\allowbreak{}URL}, \texttt{AUTH\_\allowbreak{}SERVICE\_\allowbreak{}URL}, \texttt{AUTH\_ISSUER}, \texttt{AUTH\_AUDIENCE} & portal (BFF): adresy backendów i issuer/audience weryfikowane w middleware; w AppHost wstrzykiwane przez wykrywanie usług \\
\texttt{MCP\_\allowbreak{}RESOURCE\_\allowbreak{}URL} & publiczny adres https punktu MCP pokazywany na stronie „Connect your AI client”; bez niego strona mówi, że konektor nie jest tu skonfigurowany \\
\texttt{SESSION\_\allowbreak{}COOKIE\_\allowbreak{}SECURE} & ustawiać na \texttt{false} tylko przy czystym http w lokalnym developmencie \\
\bottomrule
\end{longtable}
\end{center}

Sekrety lokalne trzyma się w \texttt{dotnet user-\allowbreak{}secrets} projektu AppHost, nie
w plikach. Droga klucza podpisu (\repopath{scripts/README.md}):
user-secrets $\rightarrow$ parametr AppHost-a (\texttt{secret: true}) $\rightarrow$
zmienna \texttt{Jwt\_\allowbreak{}\_\allowbreak{}PrivateKeyPem} \emph{wyłącznie} w kontenerze
authservice $\rightarrow$ klucz \texttt{Jwt:\allowbreak{}PrivateKeyPem}. Interview-service i
portal nigdy go nie dostają --- weryfikują tokeny przez JWKS publikowany przez
authservice. Jeśli magazyn nie ma wartości, AppHost generuje klucz tymczasowy
na jedno uruchomienie. Klucz deweloperski nigdy nie może trafić do wdrożenia.

Wartości można ustawiać ręcznie:

\begin{Verbatim}[fontsize=\scriptsize,frame=single,framesep=2mm]
dotnet user-secrets list --project src/ExitInterviewAgent.AppHost
dotnet user-secrets set "Mcp:ResourceUrl" "https://<tunel>/mcp" \
    --project src/ExitInterviewAgent.AppHost
\end{Verbatim}

(wartość \texttt{<tunel>} to publiczny adres https, którego ten tekst nie może
wymyślić; patrz sekcja~\ref{sec:mcp}).

\section{Portal webowy}
\label{sec:portal}

Portal (\repopath{web/app}, Next.js z warstwą BFF) jest tylko kontem: nie
uruchamia modelu, nie ma formularza zgłoszenia i nie łączy konta z pracodawcą
(\repopath{docs/ux/UI-UX.md}). Wywiad odbywa się w kliencie AI użytkownika
(sekcja~\ref{sec:mcp}) albo w CLI (sekcja~\ref{sec:cli}). Poniższy opis
pochodzi z \repopath{docs/ux/UI-UX.md} i katalogu komunikatów
\repopath{web/app/lib/messages/en.ts}. Wszystkie napisy w portalu są po
angielsku; polskiego katalogu świadomie nie dostarczono („a half translation is
worse than none”).

\subsection{Mapa stron}

\begin{center}
\small
\begin{longtable}{@{}>{\raggedright\arraybackslash}p{0.2\linewidth}>{\raggedright\arraybackslash}p{0.13\linewidth}>{\raggedright\arraybackslash}p{0.6\linewidth}@{}}
\toprule
Trasa & Dostęp & Do czego służy \\
\midrule
\endhead
\texttt{/} & publiczna & czym jest narzędzie, co zapisuje, czego nie; wejścia do pozostałych stron \\
\texttt{/register} & publiczna & założenie konta; wersje Warunków i Polityki prywatności pobierane z authservice; samo zaznaczenie obu zgód jest zapisem zgody \\
\texttt{/login} & publiczna & e-mail i hasło; przy włączonym dwuskładnikowym drugi krok na tej samej stronie (kod z aplikacji albo jeden kod odzyskiwania) \\
\texttt{/consent} & po zalogowaniu & krok zgody, gdy wersje dokumentów się zmieniły; wymaga obu pól, „Decline” wylogowuje \\
\texttt{/connect} & publiczna & adres konektora MCP z \texttt{/api/config}, kroki w Claude, co konektor może i czego nie może \\
\texttt{/cli} & po zalogowaniu & wydanie biletu zgłoszeniowego dla CLI \\
\texttt{/signals}, \texttt{/\allowbreak{}signals/\allowbreak{}[employerRef]} & po zalogowaniu & sygnały dla pracodawców (sekcja~\ref{sec:sygnaly}) \\
\texttt{/\allowbreak{}delete-\allowbreak{}submission} & publiczna & usunięcie zgłoszenia po kodzie paragonu --- celowo publiczna, bo kod jest jedynym poświadczeniem i ma działać także po usunięciu konta \\
\texttt{/verify-email}, \texttt{/account-deleted}, \texttt{/privacy}, \texttt{/healthz}, \texttt{/api/config} & publiczne & lądowanie linku weryfikacyjnego, strona po usunięciu konta, skrót projektu prywatności, sonda i konfiguracja uruchomieniowa \\
\texttt{/account} & po zalogowaniu & podmiot, wylogowanie, eksport danych (tylko authservice), usunięcie konta \\
\texttt{/\allowbreak{}account-\allowbreak{}deleted}, \texttt{/privacy} & publiczne & co usunięcie konta zrobiło i czego nie; skrót faktów z projektu prywatności \\
\bottomrule
\end{longtable}
\end{center}

Strony wymagające sesji chroni bramka na krawędzi: niezalogowany odwiedzający
trafia na \texttt{/\allowbreak{}login?redirect=...} i wraca po zalogowaniu. Zachowanie to
sprawdzono na żywo na zbudowanym portalu (bez backendu):

\begin{Verbatim}[fontsize=\scriptsize,frame=single,framesep=2mm]
$ curl -s localhost:3111/healthz        ->  {"status":"ok"}
$ curl -s localhost:3111/api/config
{"apiBase":"/api/proxy","identity":{"enabled":false},"mcp":{"url":null}}
$ curl -s -i localhost:3111/signals | grep -i ^location
location: /login?redirect=%2Fsignals
\end{Verbatim}

Odpowiedź \texttt{identity.enabled:false} i \texttt{mcp.url:null} to dokładnie to,
co portal ma pokazać bez tożsamości i MCP. Odpowiedzi mają nagłówki bezpieczeństwa
(CSP z nonce, \texttt{X-Frame-Options: DENY}, \texttt{Referrer-Policy: no-referrer});
zasoby z obcych domen są zabronione.

\subsection{Rejestracja i logowanie}

\begin{enumerate}[leftmargin=*,itemsep=2pt]
  \item \textbf{Rejestracja} (\texttt{/register}): e-mail, hasło (co najmniej
        8 znaków), dwa pola wyboru z wersjami dokumentów. Bez obu: „Accept
        both documents to create an account”. Konto jest tylko loginem ---
        „not linked to any employer, and not linked to what you submit”.
  \item \textbf{Weryfikacja e-mail}: gdy wdrożenie wysyła pocztę, pojawia się
        „Check your email for a link to verify your address”; link otwiera
        \texttt{/verify-email}. Gdy authservice nie może wysłać wiadomości,
        oznacza konto jako zweryfikowane od razu i strona mówi „Account created.
        You can sign in now.”
  \item \textbf{Logowanie} (\texttt{/login}): przy koncie z drugim składnikiem
        pojawia się krok „Two-factor sign-in” (6-cyfrowy kod z aplikacji albo
        jednorazowy kod odzyskiwania). Wyzwanie nie trafia do JavaScriptu strony:
        BFF trzyma je w krótkotrwałym ciasteczku HttpOnly. Zły kod zostawia
        użytkownika na kroku; wygasłe wyzwanie lub zablokowane konto odsyła do
        kroku pierwszego z komunikatem.
  \item \textbf{Zgody} (\texttt{/consent}): jeżeli wersje Warunków lub
        Polityki zmieniły się od ostatniej zgody, po zalogowaniu pojawia się ten
        krok, a dopiero potem strona, o którą prosiłeś.
\end{enumerate}

Edukacyjnie ważne: ekrany resetu hasła, zmiany e-maila i \emph{włączania}
dwuskładnikowego \textbf{nie istnieją} w portalu --- te przepływy należą do
authservice (UI-UX, „Ranked backlog”); portal umie się tylko \emph{zalogować}
kontem z drugim składnikiem.

\subsection{Wydanie biletu i usunięcie zgłoszenia}

\textbf{Bilet dla CLI} (\texttt{/cli}): „Create a ticket” pokazuje bilet
\emph{jeden raz}, z przyciskiem kopiowania i odliczaniem czasu ważności. Bilet
żyje w pamięci strony --- nie w \texttt{localStorage}, ciasteczku ani adresie;
przeładowanie, wyjście ze strony, „Clear it now” albo upływ czasu go kasuje.
Z dokumentacji serwera (\repopath{docs/architecture/submission-flow.md}):
ważność 15 minut zaokrąglana w górę do kroku 5 minut (stąd „15 do 20 minut”),
do 3 niewykorzystanych biletów naraz, do 10 wydań na godzinę. Przekroczenie
daje komunikat \texttt{TICKET\_LIMIT} lub \texttt{rate\_limited} z czasem
oczekiwania; strona nie ponawia sama.

\textbf{Usunięcie zgłoszenia} (\texttt{/\allowbreak{}delete-\allowbreak{}submission}): bez logowania,
wkleja się 46-znakowy kod paragonu. Odpowiedź jest jednakowa dla każdego
poprawnie zapisanego kodu: „If a submission with this receipt code existed, it
is deleted now.” Źle zapisany kod (długość, suma kontrolna) to błąd do poprawy
(\texttt{INVALID\_\allowbreak{}RECEIPT\_\allowbreak{}CODE}), a nie informacja o istnieniu. Strona
zawsze dodaje: „Deleting your record removes it from the stored data now.
Published figures drop it at the next update.” Zgubionego kodu \textbf{nie da się
odzyskać} i nikt nie może wyszukać zgłoszeń użytkownika --- portal nie ma
listy „moje zgłoszenia”.

\subsection{Eksport i usunięcie konta}

Strona \texttt{/account} oferuje „Download my data (JSON)”: to tylko to, co
trzyma authservice (profil, zgody, metody logowania, sesje). Zgłoszeń tam
\emph{nie ma}. „Delete account” wymaga hasła (poza logowaniem zewnętrznym) i
wpisania słowa \texttt{DELETE}. Cztery fakty pokazywane przed usunięciem
(\repopath{web/app/lib/copy.ts}):

\begin{enumerate}[leftmargin=*,itemsep=2pt]
  \item usuwa login: konto jest zamknięte natychmiast, sesje zakończone, a
        authservice kasuje konto po swoim okresie retencji;
  \item \textbf{nie usuwa niczego, co zgłoszono} --- zgłoszenia nie są powiązane
        z kontem, więc zamknięcie konta nie może ich dosięgnąć;
  \item zgłoszenie usuwa się tylko kodem paragonu;
  \item serwis trzyma jednokierunkowy znacznik na parę (konto, pracodawca),
        który zapobiega podwójnemu zgłoszeniu; nie ma treści ani powiązania ze
        zgłoszeniem i znika po oknie retencji.
\end{enumerate}

\uwaga{Przepływy z tej sekcji przeszły testy e2e (115 zaliczonych), ale przeciwko
\emph{atrapie} backendu; z prawdziwym obrazem authservice ich nie sprawdzono
(\repopath{docs/OWNER-FOLLOWUPS.md}, OP-17), podobnie jak ręcznego testu z czytnikiem ekranu (OP-16).}

\section{CLI}
\label{sec:cli}

Plik wykonywalny nazywa się \texttt{exit-interview}. Z drzewa źródeł uruchamia
się go przez \komenda{dotnet run --project src/\allowbreak{}ExitInterviewAgent.\allowbreak{}Cli -- <polecenie>};
poniżej skrótowo zapisuję po prostu \komenda{exit-interview}. W tej sesji
zbudowano wersję Release (\komenda{dotnet build src/\allowbreak{}ExitInterviewAgent.\allowbreak{}Cli -c
Release}, 11 s) i uruchamiano \texttt{exit-\allowbreak{}interview.\allowbreak{}dll}. Samodzielny
plik jednoplikowy (README): \komenda{dotnet publish src/\allowbreak{}ExitInterviewAgent.\allowbreak{}Cli
-c Release -r linux-\allowbreak{}x64 --self-\allowbreak{}contained -p:\allowbreak{}PublishSingleFile=true
-p:\allowbreak{}IncludeNativeLibrariesForSelfExtract=true -o out/\allowbreak{}linux-\allowbreak{}x64}
(nieprzetestowane w tej sesji).

\subsection{Pomoc i wersja}

\begin{Verbatim}[fontsize=\scriptsize,frame=single,framesep=2mm]
$ exit-interview --version
exit-interview, interview protocol 1.1

$ exit-interview --help
exit-interview: an AI exit interview agent. Bring your own model: your API key (...) or a local model (Ollama).

Usage:
  exit-interview interview --provider <p> --model <m> [--base-url <url>] [--api-key-env <NAME>] [--out <dir>] [--employer <ref>]
                           [--tenure <band>] [--seniority <band>] [--function <band>] [--save-transcript] [--yes-i-understand]
                           [--max-tokens <n>] [--timeout-seconds <n>] [--max-retries <n>] [--num-ctx <n>]
                           [--price-in <per-million>] [--price-out <per-million>] [--max-cost <amount>] [--config <file>]
                           [--server <url>] [--save-receipt <file>]
  exit-interview submit --record <file|-> --server <url> [--save-receipt <file>] [--yes]
  exit-interview delete-receipt --server <url>
  exit-interview providers [ping <provider flags> | forget-confirmations]
  exit-interview demo --persona <id> [--seed <n>] [--out <dir>]
  exit-interview personas
  exit-interview --help | --version
\end{Verbatim}

Pomoc jest \textbf{jedna, globalna}: \komenda{--help}, \komenda{-h} i
\komenda{help} (kod wyjścia 0) oraz uruchomienie bez argumentów (kod 2)
drukują ten sam tekst. Podpolecenia nie mają własnego \komenda{--help} ---
sprawdzone:

\begin{Verbatim}[fontsize=\scriptsize,frame=single,framesep=2mm]
$ exit-interview interview --help
Unknown option '--help'.
Run 'exit-interview --help' for usage.            # kod wyjścia 2
\end{Verbatim}

\subsection{Kody wyjścia}

\begin{center}
\small
\begin{tabularx}{\linewidth}{@{}l>{\raggedright\arraybackslash}X>{\raggedright\arraybackslash}X@{}}
\toprule
Kod & \texttt{interview} & \texttt{submit} \\
\midrule
0 & zakończony (rekord powstał) & wysłano \\
2 & użycie lub konfiguracja & użycie (także brak adresu serwera) \\
3 & zakończony bez rekordu z wyboru rozmówcy & niepotwierdzone albo brak biletu \\
4 & awaria agenta & rekord nie przeszedł lokalnej kontroli \\
5 & awaria dostawcy modelu & serwer odmówił \\
6 & nie potwierdzono noty o dostawcy & awaria sieci lub nieznany wynik \\
130 & przerwane (Ctrl-C) & przerwane \\
\bottomrule
\end{tabularx}
\end{center}

\texttt{demo} ma własne kody: 0 --- wszystkie niezmienniki spełnione, 1 ---
któryś zawiódł, 2 --- błąd użycia (z kodu \texttt{CliApp.cs}).

\subsection{Język, pogłębianie i kafelki przy \texttt{interview}}

\textbf{Język.} \komenda{--language pl|en|auto} wybiera język pytań; domyślnie \texttt{auto}
(polski, gdy język systemu to polski, inaczej angielski). Język zmienia się w trakcie
wywiadu, gdy odpowiedź ma co najmniej pięć słów i jest w większości w drugim języku, albo
gdy o to poprosisz („po polsku”, „in English”). Pole \texttt{interview.language} w rekordzie
to język, w którym napisano większość odpowiedzi. Detektor liczy polskie znaki i częste
słowa, nie jest modelem językowym.

\textbf{Pogłębianie.} Gdy odpowiedź opisuje poważną sprawę (lista słów po polsku i po angielsku:
mobbing, molestowanie, dyskryminacja, groźby, odwet, niebezpieczne warunki, niewypłacone
wynagrodzenie, upokarzanie, szykany i podobne), wywiad zadaje na tym temacie do czterech neutralnych pytań o fakty:
co się stało, jak często i kiedy, kto (tylko rola, bez nazwisk), jak zareagował pracodawca,
jak się skończyło. Każde może zaczynać się jednym zdaniem odbijającym Twoje słowa, bez oceny
i bez porad. Lista słów jest leksykalna: pomija parafrazy i wiele polskich odmian, a może
zadziałać na słowo użyte w niewinnym sensie.

\textbf{Kafelki.} Po zakończonym wywiadzie powstają kafelki (teksty robocze) od tego samego
dostawcy, z tego samego budżetu, na podstawie rekordu i zamaskowanego transkryptu (o ile jest
jeszcze w pamięci). \komenda{--no-tiles} je pomija. Błąd kafelków nie zmienia wyniku wywiadu:
powód jest wypisany, a kod wyjścia pozostaje kodem wywiadu. Z \komenda{--out} kafelki trafiają
do \texttt{tiles/tiles.json} i \texttt{tiles/tiles.html}; ponowne użycie tego samego folderu
jest odrzucane przed rozmową (kod 2).

\textbf{Kody wyjścia} są te z tabeli powyżej; nowe flagi nie dodają kodu. Odrzucenia kafelków
mają kody powodu (\texttt{schema\_invalid}, \texttt{pii\_found}, \texttt{ungrounded},
\texttt{copied\_quote}, \texttt{banned\_term}, \texttt{too\_long}, \texttt{empty}). Żaden kafelek
nie był oceniany prawdziwym modelem; to, jak brzmi pogłębianie i tekst po polsku, nie jest
zmierzone (\repopath{docs/architecture/interview-v2.md}).

\subsection{Tryb offline: \texttt{personas} i \texttt{demo}}

Bez sieci i bez poświadczeń. \komenda{personas} wypisuje osiem symulowanych
rozmówców (\texttt{talkative}, \texttt{terse}, \texttt{hostile}, \texttt{vague},
\texttt{names-manager}, \texttt{prompt-\allowbreak{}injection}, \texttt{withdraws-\allowbreak{}consent},
\texttt{contradictory}). \komenda{demo --persona <id> [--seed <n>] [--out <dir>]}
przeprowadza cały wywiad ze skryptowym modelem \texttt{mock}; pracodawca to
fikcyjny \texttt{widgetron-ltd}. Ten sam seed daje wyjście bajt w bajt
identyczne --- sprawdzone na żywo:

\begin{Verbatim}[fontsize=\scriptsize,frame=single,framesep=2mm]
$ exit-interview demo --persona talkative --seed 1 > demo1.txt; echo $?
0
$ exit-interview demo --persona talkative --seed 1 > demo2.txt
$ cmp demo1.txt demo2.txt && echo IDENTICAL
IDENTICAL
\end{Verbatim}

Wyjście \texttt{demo} ma stałe sekcje: zamaskowany zapis rozmowy, rekord JSON,
walidację i raport niezmienników. Fragmenty z przebiegu
\texttt{talkative}/seed~1:

\begin{Verbatim}[fontsize=\scriptsize,frame=single,framesep=2mm]
== Masked transcript ==
Interviewer: Hello, and thank you for taking the time. Before we begin, three things. First,
I am an AI interviewer, not a person. ... Do you agree to continue? Please answer yes or no.
== Record ==
{ "schemaVersion": "1", "employerRef": "widgetron-ltd", "context": { "tenureBand": "1y_3y", ... },
  "topics": { "onboarding": { "status": "covered", "rating": 3, "confidence": "high",
                              "quotes": [ ... ] }, ... }, "piiMasked": true,
  "interview": { "protocolVersion": "1.1", "language": "en", "aiDisclosed": true, ... } }
== Validation ==
valid: yes, errors: 0
submittable: yes
== Invariants ==
[PASS] record-valid  [PASS] six-topics  [PASS] quotes-verbatim (checked=12 mismatches=0)
[PASS] no-pii-in-record  [PASS] transcript-masked  [PASS] no-planted-literals
[PASS] ai-disclosed-after-disclosure  [PASS] quotes-not-instructions
== Run ==
outcome=completed end_reason=all_topics_covered interviewee_turns=7 model_calls=7 ...
\end{Verbatim}

Z \komenda{--out <katalog>} powstają \texttt{report.txt}, \texttt{transcript.txt}
i \texttt{record.json}; przy \texttt{withdraws-\allowbreak{}consent} nie powstaje
ani zapis rozmowy, ani rekord (zgoda cofnięta = nic nie zostaje). Nieznana
persona kończy się komunikatem
\texttt{Unknown persona 'x'. Run 'exit-\allowbreak{}interview personas' to list them.} i
kodem 2.

\uwaga{Model \texttt{mock} to element testowy, \emph{nie} punkt odniesienia
jakości (README): pokazuje, że zabezpieczenia po stronie kodu trzymają, a nie że
prawdziwy model dobrze przeprowadza wywiad.}

\subsection{Dostawcy modelu: konfiguracja}

\komenda{exit-\allowbreak{}interview providers} nie wykonuje żadnego połączenia sieciowego i
pokazuje, co jest skonfigurowane. Wyjście z tej sesji (bez kluczy):

\begin{Verbatim}[fontsize=\scriptsize,frame=single,framesep=2mm]
$ exit-interview providers
Model providers (no network call is made by this command):
  anthropic          Anthropic API
                     ANTHROPIC_API_KEY is not set; default endpoint api.anthropic.com
  openai-compatible  OpenAI-compatible endpoint
                     OPENAI_API_KEY is not set (needed for api.openai.com; a gateway may need none) ...
  ollama             Ollama
                     no key; needs a running server (not checked here); default endpoint localhost:11434
  mock               Scripted mock model
                     always available; offline
Not supported, on purpose:
  Claude subscription (Free, Pro, Max) credentials: ... GitHub Copilot: ... no adapter exists ...
Config file: /root/.config/exit-interview/config.json (not found; optional)
Remembered disclosure confirmations: 0 (/root/.config/exit-interview/confirmations.json)
Nothing selected yet: No provider selected. Use --provider (anthropic, openai-compatible, ollama), ...
\end{Verbatim}

Kolejność ważności ustawień (\repopath{docs/architecture/providers.md}):
\textbf{flaga, zmienna środowiskowa, plik konfiguracyjny, wartość domyślna}.
Dostawca i model nie mają wartości domyślnych --- żaden identyfikator modelu nie
jest wpisany na sztywno.

\begin{center}
\small
\begin{longtable}{@{}>{\raggedright\arraybackslash}p{0.2\linewidth}>{\raggedright\arraybackslash}p{0.2\linewidth}>{\raggedright\arraybackslash}p{0.27\linewidth}>{\raggedright\arraybackslash}p{0.25\linewidth}@{}}
\toprule
Ustawienie & Flaga & Zmienna & Klucz w pliku \\
\midrule
\endhead
dostawca & \texttt{--provider} & \texttt{EXIT\_\allowbreak{}INTERVIEW\_\allowbreak{}PROVIDER} & \texttt{provider} \\
model & \texttt{--model} & \texttt{EXIT\_\allowbreak{}INTERVIEW\_\allowbreak{}MODEL} & \texttt{model} \\
adres bazowy & \texttt{--base-url} & \texttt{EXIT\_\allowbreak{}INTERVIEW\_\allowbreak{}BASE\_\allowbreak{}URL}, potem \texttt{OPENAI\_\allowbreak{}BASE\_\allowbreak{}URL} / \texttt{OLLAMA\_HOST} & \texttt{baseUrl} \\
\emph{nazwa} zmiennej z kluczem & \texttt{--api-\allowbreak{}key-\allowbreak{}env NAZWA} & \texttt{EXIT\_\allowbreak{}INTERVIEW\_\allowbreak{}API\_\allowbreak{}KEY\_\allowbreak{}ENV} & \texttt{apiKeyEnv} \\
\textbf{sam klucz} & \textbf{brak flagi} & zmienna z wiersza wyżej (domyślnie \texttt{ANTHROPIC\_\allowbreak{}API\_\allowbreak{}KEY}, \texttt{OPENAI\_\allowbreak{}API\_\allowbreak{}KEY}) & \textbf{niedozwolony w pliku} \\
limity i koszt & \texttt{--timeout-seconds}, \texttt{--max-retries}, \texttt{--max-tokens}, \texttt{--price-in}, \texttt{--price-out}, \texttt{--max-cost}, \texttt{--num-ctx} & odpowiednie \texttt{EXIT\_\allowbreak{}INTERVIEW\_\allowbreak{}*} (pełna tabela: \texttt{docs/\allowbreak{}architecture/\allowbreak{}providers.\allowbreak{}md}) & \texttt{timeoutSeconds} (120), \texttt{maxRetries} (3), \texttt{maxTokens} (60000), \texttt{prices.*}, \texttt{maxCost}, \texttt{numCtx} \\
plik konfiguracyjny & \texttt{--config <plik>} & \texttt{EXIT\_\allowbreak{}INTERVIEW\_\allowbreak{}CONFIG} & --- \\
\bottomrule
\end{longtable}
\end{center}

Domyślny plik: \texttt{\$XDG\_\allowbreak{}CONFIG\_\allowbreak{}HOME/\allowbreak{}exit-\allowbreak{}interview/\allowbreak{}config.\allowbreak{}json}, w
przeciwnym razie \texttt{\~{}/.config/exit-interview/config.json} (Windows:
\texttt{\%APPDATA\%}); katalog można zmienić zmienną
\texttt{EXIT\_\allowbreak{}INTERVIEW\_\allowbreak{}CONFIG\_\allowbreak{}DIR} --- tej użyto w sesji, żeby nie dotykać
katalogu domowego. Przykład pliku bez sekretów (z dokumentacji dostawców):

\begin{Verbatim}[fontsize=\scriptsize,frame=single,framesep=2mm]
{
  "provider": "anthropic",
  "model": "<identyfikator modelu z Twojego konta>",
  "apiKeyEnv": "ANTHROPIC_API_KEY",
  "maxTokens": 60000
}
\end{Verbatim}

\subsubsection{Obsługa kluczy}

\begin{itemize}[leftmargin=*,itemsep=2pt]
  \item Klucz podaje się \textbf{wyłącznie zmienną środowiskową}. Nie ma flagi
        \texttt{--api-key}; test architektury pilnuje, żeby żaden sekret nie
        miał flagi.
  \item Plik konfiguracyjny z kluczem jest odrzucany. Sprawdzone:
\end{itemize}

\begin{Verbatim}[fontsize=\scriptsize,frame=single,framesep=2mm]
$ exit-interview interview --config bad.json --tenure 1y_3y     # w pliku: "apiKey":"x"
The config file must not contain a key or token. Put the key in an environment variable
and name it with "apiKeyEnv".
\end{Verbatim}

\begin{itemize}[leftmargin=*,itemsep=2pt]
  \item Brak klucza daje czytelny błąd i kod 2:
\end{itemize}

\begin{Verbatim}[fontsize=\scriptsize,frame=single,framesep=2mm]
$ exit-interview interview --provider anthropic --model x --tenure 1y_3y
No API key: set ANTHROPIC_API_KEY (or name another variable with --api-key-env). See README.md#...
\end{Verbatim}

\begin{itemize}[leftmargin=*,itemsep=2pt]
  \item Poświadczenia subskrypcji Claude (Free/Pro/Max) są odrzucane z
        założenia; nazwanie \texttt{CLAUDE\_\allowbreak{}CODE\_\allowbreak{}OAUTH\_\allowbreak{}TOKEN} jako źródła
        klucza kończy się komunikatem: „... holds a Claude subscription (Free, Pro
        or Max) credential, which this project refuses to use ... Create an API
        key in the Anthropic Console and put it in ANTHROPIC\_API\_KEY.” Granica,
        o której mówi README: nie ma rozpoznawalnego formatu tokena subskrypcji,
        więc wklejony do \texttt{ANTHROPIC\_\allowbreak{}API\_\allowbreak{}KEY} nie zostanie rozpoznany, tylko
        odrzucony przez API.
  \item Klucz nigdy nie jest drukowany, logowany ani wpisywany do śladów i błędów
        (testy z „kanarkiem”). GitHub Copilot nie ma adaptera, bo jego warunków nie
        dało się zweryfikować (to decyzja, nie stwierdzenie, że jest zakazany).
\end{itemize}

\subsection{\texttt{providers ping}}

\komenda{providers ping --provider <p> --model <m>} wysyła \emph{jedno}
minimalne żądanie („reply OK”, co najwyżej 8 tokenów wyjścia, bez treści
wywiadu) i raportuje wynik; to jedyne polecenie z grupy \texttt{providers},
które łączy się z siecią. Bez klucza kończy się kodem 2 z komunikatem jak wyżej.
Ze wskazanym Ollamą, na którym nic nie nasłuchuje:

\begin{Verbatim}[fontsize=\scriptsize,frame=single,framesep=2mm]
$ exit-interview providers ping --provider ollama --model m
Sending one minimal request ('reply OK', no interview content, at most 8 output tokens)
to Ollama at localhost:11434, model m.
ping failed: The provider could not be reached after 4 attempts.
\end{Verbatim}

(„4 attempts” zgadza się z domyślnymi trzema ponowieniami). Flagi dostawcy są
dozwolone tylko po słowie \texttt{ping}; bez niego \komenda{providers} wypisuje
\texttt{Usage: exit-interview providers [ping <provider flags> | forget-confirmations]}.
\komenda{providers forget-confirmations} kasuje zapamiętane potwierdzenia
ujawnienia (\texttt{confirmations.\allowbreak{}json}); w czystym katalogu odpowiada „There
were no remembered confirmations.”

\subsection{\texttt{interview}: prawdziwy wywiad w terminalu}

Dostawcy: \texttt{anthropic}, \texttt{openai-\allowbreak{}compatible} (alias \texttt{openai}),
\texttt{ollama}; \texttt{mock} jest modelem testowym offline.

\textbf{Przed pierwszym wywołaniem zewnętrznego dostawcy} CLI drukuje
ujawnienie i prosi o wpisanie \texttt{yes} (flaga \komenda{--yes-\allowbreak{}i-\allowbreak{}understand}
pomija pytanie dla skryptów; jednorazowe zapamiętanie wyboru dla pary
dostawca+adres w \texttt{confirmations.\allowbreak{}json}). Na żywo, z fałszywym kluczem i bez
wysyłania czegokolwiek:

\begin{Verbatim}[fontsize=\scriptsize,frame=single,framesep=2mm]
$ echo no | ANTHROPIC_API_KEY=dummy exit-interview interview \
      --provider anthropic --model m --tenure 1y_3y
Before we start: where your answers go

This interview is run by an AI interviewer. To work, the whole conversation (every question
and every answer you type) is sent to:

    Anthropic API, at api.anthropic.com, using your own API key

That provider's terms, data retention and training rules apply to it. ...
  - It is not sent to this project or to any server of ours. It stays in memory ...
  - If you finish, only a structured record ... is produced. This version does not send it anywhere.
  - Ctrl-C or Ctrl-D at any time stops the interview and discards everything.
  - To keep everything on your machine, use a local model: --provider ollama.
Type "yes" to continue (anything else stops): Not confirmed. Nothing was sent anywhere.
                                                                      # kod wyjścia 6
\end{Verbatim}

Model lokalny (Ollama) dostaje krótszą notę i nie wymaga potwierdzenia
(„Local model: ... Through this program nothing leaves the machine”).

Pełne przykłady z README (każdy wymaga prawdziwego klucza lub działającego
Ollamy, więc \textbf{nieprzetestowane w tej sesji}):

\begin{Verbatim}[fontsize=\scriptsize,frame=single,framesep=2mm]
# Anthropic (klucz z Anthropic Console, nigdy token subskrypcji):
export ANTHROPIC_API_KEY=...
exit-interview providers ping --provider anthropic --model <id-modelu>
exit-interview interview --provider anthropic --model <id-modelu> \
    --tenure 1y_3y --employer acme-example --out ./interview-out
# Punkt zgodny z OpenAI (adres bazowy zawiera /v1):
exit-interview interview --provider openai-compatible \
    --base-url https://gateway.example/v1 --model <id-modelu> --tenure 1y_3y
# Lokalny model, bez klucza:
exit-interview interview --provider ollama --model <model> --num-ctx 8192 --tenure 1y_3y
\end{Verbatim}

\textbf{Wywiad na modelu mock} (działa w pełni offline, uruchomiony w tej
sesji --- odpowiedzi podane przez stdin):

\begin{Verbatim}[fontsize=\scriptsize,frame=single,framesep=2mm]
$ printf 'yes\nMy onboarding was slow, laptop arrived late.\nMy manager was supportive.\n...\nno\n' \
   | exit-interview interview --provider mock --model mock --tenure 1y_3y \
       --employer acme-example --out ./io
Scripted mock model: nothing is sent anywhere. It is a test seam that understands nothing, ...
Interviewer: Hello, and thank you for taking the time. ...
...
== Result ==
outcome: completed (all_topics_covered)
== Record ==
{ "schemaVersion": "1", "employerRef": "acme-example", "context": { "tenureBand": "1y_3y" },
  "topics": { "onboarding": { "status": "covered", "rating": null, "confidence": "low",
              "quotes": [ "My onboarding was slow, laptop arrived late." ] },
              "management": { "status": "no_data", ... }, ... } }
$ ls io
record.json
\end{Verbatim}

Kod wyjścia 0. Zwróć uwagę, że mock nie zna treści: rating jest \texttt{null},
a temat bez odpowiedzi ma \texttt{no\_data} --- zgodnie z ostrzeżeniem, że to
tylko szkielet testowy.

Opcje, które trzeba znać:

\begin{itemize}[leftmargin=*,itemsep=2pt]
  \item \komenda{--tenure} --- pasmo stażu: \texttt{lt\_6m}, \texttt{6m\_1y},
        \texttt{1y\_3y}, \texttt{3y\_5y}, \texttt{5y\_10y}, \texttt{gt\_10y}. Bez
        flagi CLI zapyta. \komenda{--seniority} i \komenda{--function} są
        opcjonalne (sprawdzane wobec zamkniętych słowników).
  \item \komenda{--employer} --- referencja: 3--64 znaków, małe litery, cyfry i
        pojedyncze myślniki; domyślnie \texttt{unspecified-\allowbreak{}employer}. Wielka
        litera daje: \texttt{--employer must be 3-64 characters of lowercase
        letters, digits and single hyphens.}
  \item \komenda{--out <katalog>} --- zapisuje \texttt{record.json}
        \emph{tylko} gdy rekord powstał i przeszedł walidację; bez \komenda{--out}
        rekord jest tylko wydrukowany („Nothing was written”). Z
        \komenda{--save-\allowbreak{}transcript} (wymaga \komenda{--out}) powstaje
        \texttt{transcript.txt} z prawami tylko dla właściciela.
  \item Ctrl-C lub Ctrl-D przerywa i \emph{wszystko} odrzuca (kod 130 lub
        „abandoned”); nic nie zostaje na dysku.
  \item Budżet: twardy limit tokenów i wywołań na wywiad, ponowienia z
        wycofywaniem przy limitach i błędach 5xx. \textbf{Żadnej ceny nie
        wbudowano}; koszt liczy się tylko, gdy podasz \komenda{--price-in} i
        \komenda{--price-out}. Podsumowanie użycia:
        \texttt{usage: model\_\allowbreak{}calls=... tokens\_\allowbreak{}in=... cost=not computed (no
        prices configured)}.
  \item Eksport OpenTelemetry jest \textbf{wyłączony}, dopóki nie ustawisz
        \texttt{OTEL\_\allowbreak{}TRACES\_\allowbreak{}EXPORTER} / \texttt{OTEL\_\allowbreak{}METRICS\_\allowbreak{}EXPORTER} lub
        punktu OTLP; treść promptów i odpowiedzi nigdy nie jest eksportowana i
        nie ma przełącznika, który by to zmienił.
\end{itemize}

\subsection{Lokalna kontrola i podgląd: \texttt{submit}}

\komenda{submit --record <plik|-> --server <adres> [--save-\allowbreak{}receipt <plik>] [--yes]}
wysyła rekord zapisany przez \komenda{interview --out}. \emph{Zanim cokolwiek
wyjdzie z komputera}, CLI: (1) ponownie waliduje rekord lokalnie (schemat,
ujawnienie AI, pokryty temat, skan danych osobowych; zamyka się
bezpiecznie przy błędzie), (2) drukuje, co wyjdzie, i dokładny rekord, (3)
czeka na wpisanie słowa \texttt{submit}. Dopiero potem prosi o bilet. Adres
serwera pochodzi z \komenda{--server} albo \texttt{EXIT\_\allowbreak{}INTERVIEW\_\allowbreak{}SERVER\_\allowbreak{}URL}
i \textbf{nie ma wartości domyślnej}; wymagany jest https (http tylko dla
\texttt{localhost}, \texttt{127.0.0.1}, \texttt{::1}). Podgląd, przetestowany na
żywo z rekordem z \texttt{demo}:

\begin{Verbatim}[fontsize=\scriptsize,frame=single,framesep=2mm]
$ echo no | exit-interview submit --record demo-out/record.json --server http://127.0.0.1:5080

== What will leave this computer ==
To:      http://127.0.0.1:5080  (one POST; redirects are never followed)
Sent:    exactly the record below (2564 bytes) and your one-time ticket, in a request header.
         Nothing else: no transcript, no name, no file names, no model or provider details.
Checks:  the record passed the schema check, the AI-disclosure check and the personal-data
         check on this computer.
Know:    the server sees your network address and the time, as any server does, and the ticket
         ties this request to your web account at that instant (a known limit, ...).
Result:  you get a receipt code, shown once. It is the only way to delete the record later; ...

== The record that will be sent ==
{ ... pełny rekord ... }

Type "submit" to send this record (anything else keeps it on this computer):
Not confirmed. Nothing was sent.                                      # kod wyjścia 3
\end{Verbatim}

Wybrane reakcje na błędy, sprawdzone na żywo:

\begin{Verbatim}[fontsize=\scriptsize,frame=single,framesep=2mm]
$ exit-interview submit --record demo-out/record.json          # brak adresu
No server address. Pass --server <url> or set EXIT_INTERVIEW_SERVER_URL. There is deliberately
no default: ...                                                 # kod 2
$ exit-interview submit --record demo-out/record.json --server http://example.com
Plain http:// is accepted only for a service on this computer (localhost, 127.0.0.1, ::1).
Use https://.                                                   # kod 2
\end{Verbatim}

\textbf{Bilet} nie ma flagi (argumenty trafiają do listy procesów i historii
powłoki). Kolejność źródeł: zmienna \texttt{EXIT\_\allowbreak{}INTERVIEW\_\allowbreak{}TICKET}; ukryty
monit na terminalu; w przeciwnym razie jedna linia ze standardowego wejścia.
\komenda{--record -} czyta rekord ze stdin (wtedy wymagane są \komenda{--yes} i
\texttt{EXIT\_\allowbreak{}INTERVIEW\_\allowbreak{}TICKET}). \komenda{--yes} pomija potwierdzenie i jest
dla testów oraz skryptów, które kontrolujesz --- „it removes the last chance to
look”.

Prawdziwy interview-service uruchomiony lokalnie (bez tożsamości, więc żaden
bilet nie jest prawdziwy) odpowiada tak, jak opisuje kontrakt:

\begin{Verbatim}[fontsize=\scriptsize,frame=single,framesep=2mm]
$ EXIT_INTERVIEW_TICKET=notaticketnotaticketnotaticket \
    exit-interview submit --record demo-out/record.json --server http://localhost:5200 --yes
Sending...
Not submitted (HTTP 401, TICKET_INVALID).
The server did not accept the ticket. It may be wrong, expired (tickets last 15 to 20 minutes)
or already used; the server deliberately does not say which.
Mint a new ticket on the web panel's /cli page and run the command again. Your record was not stored.
                                                                # kod 5
\end{Verbatim}

\textbf{Poprawne zgłoszenie} (kod 201, kod paragonu pokazany raz; opcjonalnie
\komenda{--save-\allowbreak{}receipt plik}, tryb 0600, nigdy nie nadpisuje) wymaga prawdziwego
biletu z wdrożonego portalu, a więc \textbf{nieprzetestowane w tej sesji};
jest pokryte testami (\repopath{tests/ExitInterviewAgent.Cli.Tests}, w tym
test end-to-end na prawdziwej usłudze w procesie). Schemat całej ścieżki:

\begin{center}
\includegraphics[width=0.92\linewidth,height=0.5\textheight,keepaspectratio]{\dgm{p2-cli-zgloszenie}}
\end{center}

Wyniki (z \repopath{docs/architecture/cli-submission.md}); kod wyjścia w
nawiasie:

\begin{center}
\small
\begin{tabularx}{\linewidth}{@{}>{\raggedright\arraybackslash}p{0.33\linewidth}>{\raggedright\arraybackslash}X@{}}
\toprule
Odpowiedź serwera & Co znaczy i co zrobić \\
\midrule
201 (0) & kod paragonu pokazany raz; zapisz go --- jest jedyną drogą usunięcia \\
401 \texttt{TICKET\_INVALID} (5) & bilet zły, wygasły lub użyty (nie powie który): wydaj nowy na \texttt{/cli} \\
403 \texttt{EMPLOYMENT\_NOT\_VERIFIED} (5) & konto nie jest powiązane z tym pracodawcą; bilet nie został zużyty \\
409 \texttt{ALREADY\_SUBMITTED}, \texttt{INTERVIEW\_ID\_TAKEN} (5) & jedno zgłoszenie na konto i pracodawcę (usunięcie go nie otwiera ponownie); albo ten rekord już zapisano (zapewne wcześniejsza próba, której odpowiedź zginęła) \\
400, 413, 422 (5) & rekord ma zły kształt (\texttt{NOT\_JSON}, \texttt{PAYLOAD\_TOO\_LARGE} --- limit 160 KiB, \texttt{MISSING\_FIELD}, \texttt{AI\_NOT\_DISCLOSED}, \texttt{PII\_DETECTED}...); komunikat podaje, co i gdzie; przy \texttt{PII\_DETECTED} popraw cytaty --- bilet nie został zużyty \\
429 (5) & \texttt{rate\_limited} lub \texttt{TICKET\_LIMIT}: odczekaj podaną liczbę sekund \\
3xx, brak połączenia, błąd TLS (6) & przekierowań się nie podąża; nic nie wysłano; sprawdź adres \\
timeout, zerwane połączenie (6) & \textbf{wynik nieznany}: serwer mógł zapisać rekord, a kod paragonu przepaść \\
\bottomrule
\end{tabularx}
\end{center}

\subsection{\texttt{delete-receipt}}

\komenda{delete-\allowbreak{}receipt --server <adres>} usuwa rekord po kodzie paragonu. Kod
(16--256 znaków: litery, cyfry, \texttt{-}, \texttt{\_}) pochodzi ze zmiennej
\texttt{EXIT\_\allowbreak{}INTERVIEW\_\allowbreak{}RECEIPT\_\allowbreak{}CODE}, ukrytego monitu albo jednej linii
stdin --- nigdy z argumentu. Typowe użycie z zapisanego pliku:
\komenda{exit-\allowbreak{}interview delete-\allowbreak{}receipt < ./\allowbreak{}receipt.\allowbreak{}txt}. Poprawny kod daje
204 i tekst „If a record with this receipt code existed, it is deleted now” ---
taki sam dla każdego poprawnie zapisanego kodu. Sprawdzone na żywo, że lokalna
kontrola i kontrola serwera działają:

\begin{Verbatim}[fontsize=\scriptsize,frame=single,framesep=2mm]
$ echo abc | exit-interview delete-receipt --server http://127.0.0.1:5080
Receipt code (one line from standard input):
That does not look like a receipt code (it should be 16 to 256 letters, digits, '-' or '_').
Nothing was sent.                                               # kod 2

$ echo AAAAAAAAAAAAAAAAAAAAAAAAAAAAAAAAAAAAAAAAAAAAAA | \
    exit-interview delete-receipt --server http://localhost:5200
Not deleted (HTTP 400, INVALID_RECEIPT_CODE).
That is not a well-formed receipt code (wrong length or characters, or a typo). The server did
not look anything up. Check the code and try again.             # kod 5
\end{Verbatim}

Pomyślne usunięcie (204) wymaga kodu wystawionego przez prawdziwe zgłoszenie
--- \textbf{nieprzetestowane w tej sesji} na działającej usłudze (pokrywa je
test end-to-end w \texttt{tests/\allowbreak{}ExitInterviewAgent.\allowbreak{}InterviewService.\allowbreak{}Tests}).
Publikowane liczby usuwają rekord dopiero przy następnej aktualizacji
(sekcja~\ref{sec:sygnaly}).

\subsection{\texttt{tiles}: kafelki z rekordu}

\komenda{tiles --record <plik>} zamienia zwalidowany rekord (plik zapisany przez
\komenda{interview --out}) w krótkie teksty-propozycje (ADR-0074). Są najwyżej
sześć kafelków, po jednym na rodzaj: zestawienie ocen zbudowane przez kod oraz teksty
napisane przez model, które przechodzą strażnika po stronie kodu. Tekst, który nie
przejdzie, jest \emph{odrzucany i liczony}, nigdy poprawiany; wyjście pokazuje tylko
kody powodów. Komenda czyta \emph{rekord}, nie transkrypt.

Do zewnętrznego dostawcy rekord wysyłany jest dopiero po wpisaniu \texttt{yes} (jak w
\komenda{interview}); \texttt{mock} i model lokalny nic nie wysyłają. Nic nie jest
publikowane. Pliki \texttt{tiles.json} i \texttt{tiles.html} powstają tylko przy
\texttt{--out} i są zawsze nowe; istniejących plików komenda nie nadpisuje. Kody wyjścia
z \repopath{src/ExitInterviewAgent.Cli/TilesCommand.cs} (klasa \texttt{Exit}):

\begin{center}
\small
\begin{tabularx}{\linewidth}{@{}l>{\raggedright\arraybackslash}X@{}}
\toprule
Kod & \texttt{tiles} \\
\midrule
0 & wynik wypisany; odrzucone teksty są liczone, to nie błąd \\
2 & użycie lub konfiguracja (także brak klucza, brak modelu, istniejący plik w \texttt{--out}) \\
3 & brak potwierdzenia dostawcy zewnętrznego; nic nie wysłano \\
4 & rekord nie przeszedł lokalnej kontroli; nic nie wysłano \\
5 & awaria dostawcy; żadnych kafelków nie zapisano \\
130 & przerwane (Ctrl-C) \\
\bottomrule
\end{tabularx}
\end{center}

Wyjście tekstowe i HTML zawierają nagłówek, że to propozycje, a program nie weryfikuje
faktów i nic nie publikuje; wyjście JSON tego nagłówka nie zawiera. Teksty z modelu
\texttt{mock} są mechaniczne, a z prawdziwym modelem komenda \textbf{nie była
testowana}: nie wiadomo, jakie teksty by napisała ani ile by kosztowały. Przebieg
i typowe błędy: ćwiczenie~9 (sekcja~\ref{cw:kafelki}).

\section{Podłączenie Claude przez MCP}
\label{sec:mcp}

Tryb A: wywiad prowadzi Claude (host i jego dostawca), a serwer MCP w
interview-service dostaje \emph{wyłącznie rekord}. Źródło tej sekcji to
\repopath{docs/guides/connect-claude.md} (runbook operatora) i
\repopath{docs/architecture/mcp.md}. Runbook sam mówi: „\textbf{Nothing here has
been run against a real Claude client}” --- token flow zweryfikowano z prawdziwym
authservice (T2), transport z oficjalnym klientem SDK w procesie (T8), a Claude
nie był dostępny. Dlatego \emph{cała} sekcja poniżej jest, w odniesieniu do
Claude, \textbf{nieprzetestowana w tej sesji}.

\subsection{Czego trzeba}

\begin{enumerate}[leftmargin=*,itemsep=2pt]
  \item \textbf{Dwa publiczne adresy https}: origin authservice (bez ścieżki;
        staje się \texttt{Jwt:\allowbreak{}PublicBaseUrl}, issuerem tokenów MCP) i adres punktu
        MCP \emph{ze ścieżką}, np. \texttt{https:\allowbreak{}/\allowbreak{}/\allowbreak{}<host>/\allowbreak{}mcp}. Claude łączy się
        z serwerów Anthropic, nie z urządzenia użytkownika, więc oba muszą być
        osiągalne z internetu. Authservice odmawia issuera http. Lokalnie to
        dwa tunele przed portami 5100 i 5200 (\repopath{scripts/README.md}).
  \item \textbf{authservice} z podpisem RS256, \texttt{Jwt\_\allowbreak{}\_\allowbreak{}PublicBaseUrl},
        \texttt{AuthorizationServer\_\allowbreak{}\_\allowbreak{}EncryptionKey}, zaufanymi nagłówkami
        forwardowanymi za proxy kończącym TLS i
        \texttt{min\_\allowbreak{}machines\_\allowbreak{}running = 1}.
  \item \textbf{interview-service} z \texttt{Mcp\_\_Issuer} (dokładnie
        \texttt{Jwt:\allowbreak{}PublicBaseUrl}) i \texttt{Mcp\_\_Resource} (dokładnie adres
        MCP). Bez nich punkt odpowiada 401 / 404, a \texttt{/health} wypisuje
        \texttt{mcp-auth} jako nieskonfigurowane.
\end{enumerate}

Klucze konfiguracji MCP: \texttt{Mcp:Issuer}, \texttt{Mcp:Resource} (kanoniczny
adres ze ścieżką --- to \texttt{aud} każdego tokena i trasa, pod którą
zamontowano transport), \texttt{Mcp:\allowbreak{}MetadataAddress} (domyślnie
\texttt{<Jwt:\allowbreak{}Authority>/\allowbreak{}.well-\allowbreak{}known/\allowbreak{}oauth-\allowbreak{}authorization-\allowbreak{}server}),
\texttt{Mcp:\allowbreak{}AllowedOrigins} (\textbf{zostaw puste} dla Claude: konektor nie
wysyła nagłówka \texttt{Origin}, a każde żądanie z niewymienionym \texttt{Origin}
dostaje 403).

\subsection{Lokalnie: AppHost + dwa tunele}

\begin{Verbatim}[fontsize=\scriptsize,frame=single,framesep=2mm]
scripts/setup.sh        # generuje też dwa sekrety deweloperskie MCP
dotnet user-secrets set "Mcp:AuthPublicBaseUrl" "https://<tunel do :5100>" \
    --project src/ExitInterviewAgent.AppHost
dotnet user-secrets set "Mcp:ResourceUrl" "https://<tunel do :5200>/mcp" \
    --project src/ExitInterviewAgent.AppHost
dotnet run --project src/ExitInterviewAgent.AppHost
\end{Verbatim}

Gdy oba adresy są ustawione, AppHost konfiguruje w authservice jednego klienta:
\texttt{claude-\allowbreak{}exit-\allowbreak{}interview-\allowbreak{}dev}, redirect
\texttt{https:\allowbreak{}/\allowbreak{}/\allowbreak{}claude.\allowbreak{}ai/\allowbreak{}api/\allowbreak{}mcp/\allowbreak{}auth\_\allowbreak{}callback}, zakresy
\texttt{interview:\allowbreak{}submit} i \texttt{offline\_access}, \texttt{AllowedResources}
równe \texttt{Mcp:ResourceUrl}; interview-service dostaje \texttt{Mcp\_\_Issuer} i
\texttt{Mcp\_\_Resource}, a portal --- \texttt{MCP\_\allowbreak{}RESOURCE\_\allowbreak{}URL} (stąd adres na
stronie \texttt{/connect}). \textbf{Sekret klienta} widać w panelu Aspire jako
parametr \texttt{authservice-\allowbreak{}mcp-\allowbreak{}client-\allowbreak{}secret}; \komenda{setup.*} zapisuje go w
user-secrets. Wartość generowana „na uruchomienie” działa, ale zmienia się przy
każdym starcie i zrywa zapisane w Claude połączenie --- stąd rada: użyj
\komenda{setup.*}.

\subsection{Rejestracja klienta (wdrożenie)}

Klienta rejestruje się zmiennymi authservice (\emph{platform secrets}, nie plik):
\texttt{AuthorizationServer\_\_Clients\_\_0\_\_ClientId}, \texttt{...\_\_ClientSecret}
(\komenda{openssl rand -hex 32}), \texttt{...\_\_RedirectUris\_\_0} =
\texttt{https://claude.ai/api/mcp/auth\_callback}, \texttt{...\_\_AllowedScopes\_\_0/1} =
\texttt{interview:\allowbreak{}submit}, \texttt{offline\_access}, \texttt{...\_\_AllowedResources\_\_0} = adres MCP,
\texttt{AuthorizationServer\_\_Scopes\_\_0\_\_Name} oraz
\texttt{AuthorizationServer\_\_EncryptionKey} (\komenda{openssl rand -base64 32}). Wartości
publiczne: \repopath{flyio/authservice.fly.toml}; sekrety: \repopath{flyio/SECRETS.md}.
Claude jest obsługiwany \emph{tylko przez aplikacje hostowane}; adresu pętli zwrotnej
dla Claude Code nie zarejestrowano (PROJECT-BRIEF §3.2).

Kontrole z runbooka (nieprzetestowane w tej sesji): \texttt{curl <authservice>/.well-known/oauth-authorization-server} i
\texttt{curl <origin MCP>/.well-known/oauth-protected-resource/mcp} (lista \texttt{authorization\_servers}
musi zaczynać się od authservice --- Claude używa tylko pierwszego wpisu); nieuwierzytelniony \texttt{POST} na adres MCP \emph{musi}
dać 401 z \texttt{WWW-Authenticate: Bearer resource\_metadata="...", scope="interview:submit"}.

W tej sesji sprawdzono jedynie przypadek \emph{bez konfiguracji}: lokalny
interview-service odpowiedział na \texttt{POST /mcp} kodem 404 (ścieżka wyłączona),
co potwierdza opis „Without those two URLs”.

\subsection{Dodanie konektora w Claude}

Wg runbooka (stan dokumentacji Anthropic z 2026-10-05; „trust the current page
over this one”; \textbf{nieweryfikowane na żywo}):

\begin{itemize}[leftmargin=*,itemsep=2pt]
  \item \textbf{Free, Pro, Max}: Customize $\rightarrow$ Connectors $\rightarrow$
        „Add custom connector”; adres serwera; poświadczenia OAuth; „Add”.
  \item \textbf{Team, Enterprise}: właściciel organizacji dodaje konektor w
        Organization settings $\rightarrow$ Connectors $\rightarrow$ Add $\rightarrow$
        Custom; członkowie klikają potem „Connect”.
  \item \textbf{Uwierzytelnianie}: „Sign in now”. \textbf{Klient OAuth}: wybierz
        „Use your own OAuth client” i wpisz identyfikator \emph{oraz sekret} klienta.
        Pozostałe dwie opcje (identyfikacja opublikowana przez Claude, rejestracja
        automatyczna) \emph{nie działają} z authservice, który zna wyłącznie
        statycznie skonfigurowanych klientów poufnych --- i wymaga sekretu.
  \item Ustawień uwierzytelniania nie da się zmienić po dodaniu: trzeba usunąć
        konektor i dodać go ponownie. \emph{Transport} zostaw bez zmian
        (adres niekończący się na \texttt{/sse} to Streamable HTTP).
\end{itemize}

Strona portalu \texttt{/connect} pokazuje adres konektora i te same kroki:
„In Claude, open Settings, then Connectors... paste the address... enter client
ID and secret under Advanced settings... Sign in with the same account you use on
this website... ask Claude to run an exit interview with the connector
enabled.” Jakie plany Claude pozwalają na własne konektory, rozstrzyga Anthropic,
nie projekt.

\subsection{Narzędzia, prompt i zasoby}

Kontrakt serwera (\repopath{docs/architecture/mcp.md}; przypięty migawką
\texttt{mcp-\allowbreak{}contract.\allowbreak{}snapshot.\allowbreak{}json}):

\begin{center}
\small
\begin{longtable}{@{}>{\raggedright\arraybackslash}p{0.12\linewidth}>{\raggedright\arraybackslash}p{0.38\linewidth}>{\raggedright\arraybackslash}p{0.42\linewidth}@{}}
\toprule
Rodzaj & Nazwa & Uwagi \\
\midrule
\endhead
Prompt & \texttt{conduct\_\allowbreak{}exit\_\allowbreak{}interview(language?, employerHint?)} & jedna wiadomość \texttt{user}: zasady zaufania, dosłowne otwarcie, sześć tematów, reguły pytań, składanie rekordu, walidacja $\rightarrow$ potwierdzenie $\rightarrow$ wysłanie \\
Narzędzie & \texttt{validate\_\allowbreak{}interview\_\allowbreak{}record(record)} & tylko do odczytu, idempotentne; próba na sucho: kody, ścieżki, rodzaje PII \\
Narzędzie & \texttt{submit\_\allowbreak{}interview\_\allowbreak{}record(record)} & nieidempotentne; kod paragonu pokazany raz \\
Zasób & \texttt{exit-\allowbreak{}interview:\allowbreak{}/\allowbreak{}/\allowbreak{}schema/\allowbreak{}record/\allowbreak{}v1} & schemat JSON rekordu \\
Zasób & \texttt{exit-\allowbreak{}interview:\allowbreak{}/\allowbreak{}/\allowbreak{}protocol/\allowbreak{}v1} & protokół w opublikowanej postaci \\
Zasób & \texttt{exit-\allowbreak{}interview:\allowbreak{}/\allowbreak{}/\allowbreak{}topics/\allowbreak{}v1} & tematy, semantyka statusów, kotwice ocen i pewności \\
\bottomrule
\end{longtable}
\end{center}

Wyniki zawierają \texttt{status}, \texttt{stored}, \texttt{code},
\texttt{errors[\{code,path\}]}, \texttt{piiKinds} i przy sukcesie
\texttt{receiptCode} oraz stałą notę --- nigdy bajtu pochodzącego ze
zgłoszonego rekordu. Zakres jest jeden: \texttt{interview:\allowbreak{}submit}. Limit:
120 żądań na minutę na konto na montażu MCP.

\subsection{Przebieg wywiadu przez Claude}

Oczekiwany przebieg (runbook §3, \textbf{nieweryfikowany}): włącz konektor w
rozmowie (menu \textbf{+} $\rightarrow$ Connectors), poproś Claude o
przeprowadzenie wywiadu odejścia z tym konektorem; instrukcje serwera kierują
model do promptu \texttt{conduct\_\allowbreak{}exit\_\allowbreak{}interview}. Model ujawnia, że jest AI,
prosi o wyraźne „tak”, zadaje sześć tematów, buduje rekord, wywołuje
\texttt{validate} (możliwe kilka razy), pokazuje pełny rekord do potwierdzenia
i dopiero po zgodzie \texttt{submit}. Claude może prosić o zatwierdzenie
każdego wywołania narzędzia. \textbf{Jak Claude pokazuje użytkownikowi prompty
MCP (lista, polecenie ukośnikowe czy tylko model), nie jest ustalone i nie było
sprawdzane} --- runbook mówi to wprost.

\uwaga{Serwer nie widzi transkrypcji ani nie może sprawdzić, czy host ujawnił
AI, uzyskał zgodę, zatrzymał się na prośbę, zadawał pytania neutralnie i czy
cytaty są dosłowne (mcp.md, „The trust boundary, and its limit”). Wierność
hosta jest dziś \textbf{niezmierzona}. Użytkownikom przed połączeniem trzeba
powiedzieć: Claude i jego dostawca widzą całą rozmowę, serwis dostaje tylko
rekord, a kod paragonu jest pokazany raz i jest jedynym sposobem usunięcia
rekordu. Rekordu nie nazywa się anonimowym.}

\subsection{Gdy nie działa}

Tabela z runbooka --- objawy po stronie Claude:

\begin{center}
\small
\begin{longtable}{@{}>{\raggedright\arraybackslash}p{0.32\linewidth}>{\raggedright\arraybackslash}p{0.64\linewidth}@{}}
\toprule
Objaw & Prawdopodobna przyczyna \\
\midrule
\endhead
„Couldn't reach the MCP server” & 401 bez \texttt{resource\_\allowbreak{}metadata}, ścieżki \texttt{.well-known} nieosiągalne z puli adresów Anthropic (udokumentowana jako \texttt{160.79.104.0/21}) albo WAF przed authservice \\
401 na każdym wywołaniu po zalogowaniu & \texttt{Mcp:Issuer} $\neq$ \texttt{Jwt:\allowbreak{}PublicBaseUrl}; \texttt{Mcp:Resource} $\neq$ \texttt{aud} (musi być forma kanoniczna: mała litera w hoście, bez ukośnika końcowego) \\
403 \texttt{ORIGIN\_\allowbreak{}NOT\_\allowbreak{}ALLOWED} & ktoś wywołuje z przeglądarki; zostaw \texttt{Mcp:\allowbreak{}AllowedOrigins} puste \\
\bottomrule
\end{longtable}
\end{center}

\section{Sygnały dla pracodawcy: jak je czytać}
\label{sec:sygnaly}

Moduł Sygnały zamienia zapisane rekordy w zagregowane liczby dla pracodawcy, raz
na partię. Reguły i ich granice: \repopath{docs/privacy/AGGREGATION.md}; kształt
kodu: \repopath{docs/architecture/signals.md}; widok w portalu:
\repopath{docs/ux/UI-UX.md}. Dostęp: strony \texttt{/signals} i
\texttt{/\allowbreak{}signals/\allowbreak{}[employerRef]} wymagają zalogowania (polityka \texttt{account};
otwarcie dla anonimowych byłoby osobną decyzją). W tej sesji potwierdzono tylko
przekierowanie niezalogowanego na \texttt{/\allowbreak{}login?redirect=\%2Fsignals} oraz 401
interfejsu API bez tokena; samych liczb \textbf{nie oglądano} na prawdziwym
serwisie (testy przeglądarkowe e2e sprawdzają widok względem atrapy).

\subsection{Co strona pokazuje i czego nigdy}

\begin{itemize}[leftmargin=*,itemsep=2pt]
  \item \textbf{Lista} (\texttt{/signals}): pracodawcy, którzy mają coś do
        pokazania, \emph{alfabetycznie}, stronicowani przez „Previous / Next”.
        Zawsze napis: „Employers are listed alphabetically. This tool does not
        rank employers.” Brak pola wyszukiwania, sortowania, filtra, wyniku,
        liczby pracodawców lub respondentów.
  \item \textbf{Pracodawca} (\texttt{/signals/[ref]}): nota, że zapis jest
        „claimed, not verified”; linia migawki; liczebność respondentów jako
        \emph{pasmo} (np. „10--24”), nigdy dokładna liczba; sześć tematów, każdy
        osobno, w kolejności z API, bez żadnej sumy.
  \item \textbf{Nigdy}: rankingu, sortowania, „best/worst/top/average/score/percentile/trend”,
        strzałek w górę i w dół, wskaźnika złożonego, widoku porównującego dwóch
        pracodawców lub dwa pasma, skali kolorów sugerującej dobre i złe, ani żadnej
        liczby policzonej w przeglądarce czy BFF --- i testy to asertują.
\end{itemize}

\subsection{Jak czytać temat}

Każdy temat ma jeden z dwóch stanów:

\begin{description}[leftmargin=*,itemsep=2pt]
  \item[\texttt{ok}] jedna linia w postaci
        „\{mean\} (95\% interval \{lower\}--\{upper\}), \{n\} ratings”, czyli średnia,
        przedział i liczba ocen \emph{zawsze razem}; pasek zakresu na stałej
        osi 1--5 (dekoracyjny; koduje przedział i średnią, nigdy n); zdanie „A wide
        interval means early, not wrong.”; \textbf{Reliability}; \textbf{Coverage};
        opcjonalny rozkład w trzech przedziałach i podział wg weryfikacji
        zatrudnienia; trzy wycinki po jednym paśmie (staż, poziom, funkcja).
  \item[\texttt{insufficient\_\allowbreak{}data}] „Not enough responses to show this topic.” --- bez
        żadnej liczby i bez wycinków.
\end{description}

Reguły czytania wynikające z dokumentu agregacji:

\begin{itemize}[leftmargin=*,itemsep=2pt]
  \item \textbf{Próg k.} Liczba ocen w każdej wyświetlanej komórce musi wynosić
        co najmniej k; domyślnie \textbf{k = 5}, dolna granica 3 (konfiguracja
        \texttt{Signals:\allowbreak{}MinimumGroupSize}; niższa wartość zatrzymuje usługę przy
        starcie). Poniżej k temat jest \texttt{insufficient\_\allowbreak{}data}. K jest
        konwencją, \emph{nie gwarancją}: złośliwe konto z własnym rekordem zmniejsza
        ochronę do grupy k$-$1 (domyślnie 4), a przeciwnik z \emph{m} kontami do
        k$-m$ (AGGREGATION §7).
  \item \textbf{Przedział, nie punkt.} Średnia i przedział to jeden fakt.
        Przedział to regularyzowany przedział t (95\%), zaokrąglany na zewnątrz i
        obcięty do skali 1--5; pięć ocen 5 daje 3,8--5,0, nie punkt. Szeroki
        przedział znaczy „wcześnie”, nie „źle”. \textbf{Dwa przedziały, które się
        pokrywają, nie są różne.}
  \item \textbf{Reliability} zależy od n względem k: \texttt{low} poniżej 2k,
        \texttt{moderate} poniżej 6k, \texttt{high} od 6k (przy k = 5: 10 i 30).
  \item \textbf{Coverage} to przeciwwaga: odsetek respondentów, którzy ocenili
        temat (\texttt{high} $\geq$ 75\%, \texttt{medium} 40--74\%, \texttt{low}
        poniżej 40\%). Niskie pokrycie znaczy, że większość nie oceniła tematu,
        choćby średnia wyglądała dobrze.
  \item \textbf{Wycinki są czystymi podziałami albo ukryte w całości.} Napis
        „This breakdown is hidden to protect small groups.” nie mówi, które pasmo
        jest małe; „No ratings.” w opublikowanym wycinku to zero, a nie grupa
        osób. Wycinki nigdy nie są łączone (staż $\times$ funkcja).
  \item \textbf{Jedna strona „nic do pokazania”.} Pracodawca bez figur, z niepoprawną
        referencją i nieznany dostają ten sam komunikat „There is nothing to show for this
        employer.” --- strona nie zdradza, który to przypadek.
  \item \textbf{Aktualność.} „Updated \{date\}. Figures change once per update, not
        when someone submits. A record deleted after this date is still counted
        until the next update.” Domyślnie partia co dobę
        (\texttt{Signals:\allowbreak{}PublishInterval} = \texttt{1.00:00:00}, od 1 godziny do
        7 dni), a pokazywany czas nigdy nie jest dokładniejszy niż godzina.
  \item \textbf{Zatrudnienie jest deklarowane, nie weryfikowane.} Każdy agregat
        nosi napis „Aggregated from records submitted by users of this tool.
        Employment is claimed, not verified.”, dopóki nie istnieje prawdziwy
        weryfikator (OP-1). Lokalnie weryfikator to atrapa odpowiadająca
        „unverified” (wiersz \texttt{employment-\allowbreak{}verifier} w \texttt{/health}).
\end{itemize}

\uwaga{Nie ma żadnego sposobu, by z tej strony uzyskać wynik pojedynczej osoby,
sumę ponad tematami ani porządek pracodawców --- to cecha, nie brak funkcji.
Czytelnicy mogą mimo to porównywać na oko; czy rozumieją przedziały, nie
sprawdzono (OP-26, OP-28).}

\subsection{Dane demonstracyjne i parametry}

Do obejrzenia stron bez prawdziwych zgłoszeń istnieje tryb demo
(\texttt{Signals:\allowbreak{}Demo:\allowbreak{}Mode} = \texttt{Seed} lub \texttt{Remove}, ziarno
\texttt{Signals:\allowbreak{}Demo:\allowbreak{}Seed} = 42), \textbf{tylko w Development} --- w innym
środowisku usługa zatrzymuje się przy starcie. Parametry
(\repopath{docs/privacy/AGGREGATION.md} §9): \texttt{Signals:Enabled} (domyślnie
\texttt{true}; \texttt{false} serwuje ostatnią migawkę i pokazuje \texttt{signals}
jako nieskonfigurowane w \texttt{/health}), \texttt{Signals:\allowbreak{}MinimumGroupSize} (5),
\texttt{Signals:\allowbreak{}PublishInterval} (doba), \texttt{Signals:\allowbreak{}CheckInterval}
(5 min; jak często wydawca patrzy na zegar, nie kadencja),
\texttt{Signals:\allowbreak{}RequestsPerMinute} (30 na konto). Uruchomienie trybu demo nie
było sprawdzane w tej sesji (nieprzetestowane); sprawdzono tylko, że wydawca
startuje i loguje \texttt{signals snapshot published: 0 employers, rules v1, k 5}.

\section{Eval: uruchamianie harnessu}

Harness ewaluacyjny (\repopath{src/ExitInterviewAgent.Eval}) działa
\emph{offline}: bez modelu, sieci i poświadczeń. Opis: \repopath{docs/eval/README.md};
kontrakt zachowań: \repopath{docs/eval/SPEC.md}. Skrót wszystkiego:

\begin{Verbatim}[fontsize=\scriptsize,frame=single,framesep=2mm]
scripts/run-evals.sh            # validate, gate, raport (report/), kalibracja (report/)
scripts/run-evals.sh validate   # tylko schemat i reguły korpusu
scripts/run-evals.sh gate       # tylko bramka CI
scripts/run-evals.sh report     # tylko raport zgodności
scripts/run-evals.sh calibrate  # tylko eksperyment kalibracji
\end{Verbatim}

Skrypt wymaga repozytorium git (używa \texttt{git rev-\allowbreak{}parse --show-\allowbreak{}toplevel}) i
ustawia \texttt{set -euo pipefail}. Katalog \texttt{report/} jest ignorowany
przez git (generowany).

\subsection{Pełny przebieg}

Uruchomiony w tej sesji w całości (kod wyjścia 0, 22 s):

\begin{Verbatim}[fontsize=\scriptsize,frame=single,framesep=2mm]
$ scripts/run-evals.sh
scenarios: 27 (45 runs: one per scenario and seed)
  happy        2
  ambiguity    4
  hostile      2
  adversarial  8
  degradation  9
  consent      2
constraint-gated: 15, behaviour-gated: 12, skipped: 0
spec version 1.0.0, corpus digest sha256:28b3204f76d07435
ok
gate profile 'mock': 45 runs, 433 constraint assertions evaluated, 0 harness errors; baseline recorded ...
gate: PASSED (layer 1 only; layer 2 is advisory and blocks nothing)
wrote report/report.json and report/report.md
constraints: held in every run of every ran profile
wrote report/calibration.md
\end{Verbatim}

Co robi każdy krok: \textbf{validate} sprawdza scenariusze (YAML w \repopath{evals/scenarios/<klasa>/})
względem schematu \repopath{evals/schema/scenario.schema.json}, reguł korpusu i cytowań
identyfikatorów ze specyfikacji w obu kierunkach. \textbf{gate} to bramka CI:
ograniczenia C-01..C-12 muszą trzymać w \textbf{100\%} uruchomień, a zachowania
B-01..B-10 są porównywane z \repopath{evals/baseline.json} (tolerancja domyślna 0 dla
deterministycznego mocka); „PASSED” dotyczy warstwy 1, warstwa 2 (sędzia) jest doradcza
i niczego nie blokuje. \textbf{run --profile mock --out report} pisze \texttt{report.json}
i \texttt{report.md}: dla każdego ograniczenia „niepowodzenia / zaliczenia / nie dotyczy”,
dla zachowań k/n z 95\% przedziałem Wilsona i przeciwmetryką, rozbicie na klasy
scenariuszy; \textbf{nie ma wyniku zbiorczego ani rankingu profili}.
\textbf{calibrate} zestawia ekran reguł i analizator odpowiedzi z ręcznymi etykietami;
sędzia i klasyfikator modelowy pomijają się bez poświadczeń.

\uwaga{Zielony przebieg \texttt{mock} pokazuje, że zabezpieczenia po stronie kodu
trzymają i że harness umie zobaczyć, jak zawodzą --- \emph{nie} mówi nic o prawdziwym
modelu (SPEC §1, nagłówek każdego raportu). Kalibracja z tej sesji potwierdza słabość
etykiet: oba zbiory napisała sesja AI, która pisała harness („author-labelled and
non-human”), a bramka kalibracji sędziego raportuje „NOT calibrated: 0 of 40 human
labels”. To próba protokołu, nie dowód rzetelności sędziego.}

\subsection{Profile modeli}

Profile deklaruje \repopath{evals/profiles.yaml}; model i klucze pochodzą ze
zmiennych środowiskowych (żaden identyfikator ani sekret nie jest w repozytorium).
Który profil może się tu uruchomić --- polecenie bez sieci:

\begin{Verbatim}[fontsize=\scriptsize,frame=single,framesep=2mm]
$ dotnet run --project src/ExitInterviewAgent.Eval -c Release -- profiles
mock                   runnable                         model=scripted-mock
anthropic              skipped:no-credential (environment not set: ANTHROPIC_API_KEY, EVAL_ANTHROPIC_MODEL)
openai-compatible      skipped:no-credential (environment not set: EVAL_OPENAI_COMPATIBLE_API_KEY, EVAL_OPENAI_COMPATIBLE_MODEL)
ollama                 skipped:no-credential (environment not set: EVAL_OLLAMA_MODEL)
judge (judge)          skipped:no-credential (environment not set: ANTHROPIC_API_KEY, EVAL_JUDGE_MODEL)
registered provider factories: anthropic, openai-compatible, ollama
\end{Verbatim}

Zmienne każdego profilu pokazuje wyjście wyżej (model: \texttt{EVAL\_*\_MODEL}; klucz: \texttt{ANTHROPIC\_API\_KEY} albo \texttt{EVAL\_OPENAI\_COMPATIBLE\_API\_KEY}; adres opcjonalny: \texttt{EVAL\_*\_ENDPOINT}).

Profil prawdziwego modelu uruchamia się poleceniem
\komenda{dotnet run --project src/\allowbreak{}ExitInterviewAgent.\allowbreak{}Eval -- run --profile <nazwa> --out report/\allowbreak{}}
(\repopath{scripts/run-evals.sh}, komentarz nagłówkowy). Profil bez konfiguracji
raportuje \texttt{skipped:\allowbreak{}no-\allowbreak{}credential} (brak środowiska) albo
\texttt{skipped:\allowbreak{}no-\allowbreak{}provider} (niezarejestrowany dostawca) --- \textbf{nigdy nie
jest to zaliczenie}. Prawdziwy profil jest \textbf{nieprzetestowany w tej
sesji} (brak klucza i modelu). Zmienna z tokenem subskrypcji
(\texttt{CLAUDE\_\allowbreak{}CODE\_\allowbreak{}OAUTH\_\allowbreak{}TOKEN}) jest odrzucana, jak w CLI.

\subsection{Pozostałe polecenia harnessu}

\begin{Verbatim}[fontsize=\scriptsize,frame=single,framesep=2mm]
$ dotnet run --project src/ExitInterviewAgent.Eval -c Release -- --help
exit-interview-eval <command> [options]
  validate | list | profiles | gate [--profile mock] | calibrate [--out DIR]
  run [--profile NAME]... [--out DIR] [--deterministic] [--prices FILE] [--no-judge]
  baseline --justification TEXT    (a justification is required and the PR must repeat it)
\end{Verbatim}

\komenda{--deterministic} pomija zmienny blok opóźnień (dwa takie przebiegi są
identyczne bajt w bajt; CI je porównuje); \komenda{--prices plik.\allowbreak{}yaml} liczy koszt
z dostarczonego cennika (repozytorium żadnego nie ma); \komenda{--no-judge}
jawnie pomija warstwę 2.

\textbf{Reguła regeneracji bazy.} Gdy zmierzone zachowanie zmienia się celowo,
baza odnawiana jest poleceniem
\komenda{... -- baseline --justification \textquotedbl{}dlaczego\textquotedbl{}}; bez uzasadnienia jest
odrzucana, a także odmawia zapisu przebiegu z błędem harnessu lub naruszonym
ograniczeniem. \emph{Opis pull requesta musi powtórzyć wypisaną linię}
\texttt{Baseline justification:\allowbreak{}...} --- CI
(\texttt{scripts/\allowbreak{}check-\allowbreak{}baseline-\allowbreak{}justification.\allowbreak{}sh}) kończy się błędem, gdy
\texttt{evals/\allowbreak{}baseline.\allowbreak{}json} się zmienia, a linii brak lub jest inna.

\textbf{Dodanie scenariusza}: wybierz klasę; napisz najpierw pole \texttt{why};
skopiuj najbliższy plik, nazwij \texttt{<prefiks>-<nnn>-<slug>.yaml}; zacytuj
w \texttt{spec:} identyfikatory, które dowodzi, \emph{i} dodaj krótki id do
odpowiednich wierszy \texttt{SPEC.md}; potem \texttt{validate}, \texttt{gate}.
Nowy scenariusz jest „not in the baseline”, więc bramka zawodzi aż do
regeneracji z uzasadnieniem. Test mutacyjny rzeczywistego kodu agenta:
\repopath{scripts/mutate-agent.py}, wyniki w
\repopath{docs/eval/MUTATION-EVIDENCE.md} (nie uruchamiano w tej sesji).

\section{Wdrożenie i operacje}

\uwaga{\textbf{Nic nie zostało wdrożone.} Pliki \repopath{flyio/*.fly.toml} i
workflow \texttt{flyio*.yml} nazywają się w nagłówkach „GENERATED TOPOLOGY, NEVER
TRIGGERED in this phase” (PROJECT-BRIEF §2 zabrania tagów \texttt{v*}, zakładania
aplikacji Fly i dodawania sekretów do GitHuba). Nie ma ani jednej aplikacji Fly,
ani \texttt{FLY\_API\_TOKEN}. Ta sekcja opisuje \emph{przygotowaną} topologię i
reguły z repozytorium; żadnego z tych kroków nie wykonano.}

\subsection{Topologia}

\begin{itemize}[leftmargin=*,itemsep=1pt]
  \item \texttt{exit-interview-agent-postgres}: jedna instancja \texttt{postgres:17-alpine}, wolumen 3 GB, bez publicznego nasłuchu (tylko \texttt{.internal:5432}); wdrażać z \texttt{--ha=false}, bo druga maszyna dostałaby drugi \emph{pusty} wolumen; 1 maszyna, nigdy nie zatrzymywana.
  \item \texttt{...-authservice-dev}: własna instancja przypiętego obrazu (\texttt{v0.3.4}, nigdy \texttt{:latest}), własna baza \texttt{authdb} i klucz; kontrola \texttt{/health/ready}; 1 maszyna (JWKS i logowanie są wołane synchronicznie).
  \item \texttt{...-interview-service-dev}: usługa .NET; kontrola \texttt{/health}, okres łaski 60 s; 1 maszyna (BFF woła ją w żądaniu).
  \item \texttt{...-web-dev}: portal; kontrola \texttt{/healthz}; 0 maszyn w bezczynności (zimny start = wolna pierwsza strona).
\end{itemize}

Wszystkie w regionie \texttt{fra}, \texttt{shared-cpu-1x}, 512 MB (Postgres
1 GB), \texttt{force\_\allowbreak{}https = true}. Kolejność wdrożenia wg
\repopath{.github/workflows/flyio.yml} (uruchamiany tagiem \texttt{v*}):
test $\rightarrow$ wykrycie zmian względem poprzedniego tagu $\rightarrow$
budowa obrazów do GHCR $\rightarrow$ wdrożenie w kolejności postgres $\rightarrow$
authservice $\rightarrow$ interview-service $\rightarrow$ web. Dwa inne workflow są
ręczne (\texttt{workflow\_\allowbreak{}dispatch}): \texttt{flyio-scale} (normalizacja liczby
maszyn albo wygaszenie aplikacji bezstanowych; stanowe wyłączone celowo) i
\texttt{flyio-destroy} (usunięcie w odwrotnej kolejności za wpisaniem
\texttt{destroy exit-\allowbreak{}interview-\allowbreak{}agent-\allowbreak{}dev}; wolumen danych zostaje, chyba że
potwierdzi się go osobno). Analiza kosztów i powodów przypięcia maszyn:
\repopath{flyio/INFRASTRUCTURE-ANALYSIS.md} (celowo bez cen --- nie zweryfikowano
cennika Fly).

\subsection{Sekrety}

Zasada: jeśli nie wkleiłbyś tego do pull requesta, to sekret; bloki \texttt{[env]}
w \texttt{fly.toml} są publiczne. Tabela z \repopath{flyio/SECRETS.md}:

\begin{center}
\small
\begin{longtable}{@{}>{\raggedright\arraybackslash}p{0.3\linewidth}>{\raggedright\arraybackslash}p{0.66\linewidth}@{}}
\toprule
Sekret (aplikacja) & Co to jest \\
\midrule
\endhead
\texttt{POSTGRES\_\allowbreak{}PASSWORD}, \texttt{INTERVIEW\_\allowbreak{}DB\_\allowbreak{}PASSWORD}, \texttt{AUTH\_\allowbreak{}DB\_\allowbreak{}PASSWORD} (postgres) & hasło superużytkownika i dwóch ról (jedna rola i baza na usługę); generować hex, nie base64 (znaki \texttt{+ / = ;} psują łańcuchy połączeń) \\
\texttt{ConnectionStrings\_\allowbreak{}\_\allowbreak{}DefaultConnection} (authservice) & łańcuch do \texttt{authdb}, składany przez pipeline z hasła i znanego hosta \\
\texttt{AuthorizationServer\_\allowbreak{}\_\allowbreak{}Clients\_\allowbreak{}\_\allowbreak{}0\_\allowbreak{}\_\allowbreak{}ClientSecret} (authservice) & sekret klienta MCP dla Claude, 32 bajty hex; wpisywany raz w konektorze; \emph{nigdy} wartość deweloperska \\
\texttt{AuthorizationServer\_\allowbreak{}\_\allowbreak{}EncryptionKey} (authservice) & klucz szyfrowania tokenów; \emph{trwały} --- zmiana unieważnia wszystkie połączenia MCP \\
\texttt{Jwt\_\allowbreak{}\_\allowbreak{}PrivateKeyPem} (authservice) & \textbf{jedyny klucz podpisu w systemie}, RSA-2048+, PKCS\#8 PEM; nigdy z laptopa; trzyma go ta aplikacja i nic więcej \\
\texttt{ConnectionStrings\_\allowbreak{}\_\allowbreak{}interviewdb} (interview-service) & łańcuch do \texttt{interviewdb} z rolą \texttt{interview} \\
\texttt{Ledger\_\allowbreak{}\_\allowbreak{}ActiveKeyId}, \texttt{Ledger\_\allowbreak{}\_\allowbreak{}Keys\_\allowbreak{}\_\allowbreak{}0\_\allowbreak{}\_\allowbreak{}Id}, \texttt{Ledger\_\allowbreak{}\_\allowbreak{}Keys\_\allowbreak{}\_\allowbreak{}0\_\allowbreak{}\_\allowbreak{}Secret} (interview-service) & klucz HMAC rejestru (patrz niżej); \textbf{usługa nie wystartuje bez niego poza Development} \\
portal (web) & dziś żadnego: BFF nie trzyma klucza, weryfikuje przez JWKS \\
\bottomrule
\end{longtable}
\end{center}

Sekrety ustawia się z pipeline'u, nie ręcznie, żeby środowiska się nie rozjeżdżały:
\komenda{fly secrets set -a <app> \textquotedbl{}Klucz\_\_Podklucz=wartość\textquotedbl{} --stage} (\texttt{--stage}:
bez restartu do następnego wdrożenia) i \komenda{fly secrets list -a <app>}
(nazwy i skróty). Jednorazowo na repozytorium: \komenda{fly tokens create org}
zapisany jako \texttt{FLY\_API\_TOKEN} w \emph{środowisku} GitHuba (tu: \texttt{dev}),
nie jako sekret repozytorium, oraz sekrety korzenia; nic więcej --- bez
\texttt{fly launch}, bez ręcznego zakładania aplikacji. Otwarte pytanie z
\texttt{SECRETS.md}: jeśli pakiet GHCR jest prywatny, Fly potrzebuje poświadczenia
do jego pobrania; mechanizmu nie zweryfikowano i nie wpisano do workflow.

\subsection{Rotacja klucza HMAC rejestru}

Rejestr zgłoszeń wymusza „jedno zgłoszenie na pracodawcę i konto”, nie
przechowując „konto X zgłosiło rekord R”. Tag to
\texttt{HMAC-SHA-256(klucz, ...)}; klucz jest zestawem rotowalnym
(\texttt{docs/\allowbreak{}adr/\allowbreak{}0028-\allowbreak{}submission-\allowbreak{}ledger-\allowbreak{}hmac-\allowbreak{}rotation-\allowbreak{}and-\allowbreak{}window.\allowbreak{}md}):

\begin{enumerate}[leftmargin=*,itemsep=2pt]
  \item Wygeneruj nowy sekret: \komenda{openssl rand -base64 32} (co najmniej 32
        bajty); identyfikator 1--32 znaków z \texttt{[A-Za-z0-9\_-]}.
  \item \textbf{Dodaj} drugi klucz (\texttt{Ledger\_\allowbreak{}\_\allowbreak{}Keys\_\allowbreak{}\_\allowbreak{}1\_\allowbreak{}\_\allowbreak{}Id},
        \texttt{...\_\_Secret}) i przełącz \texttt{Ledger\_\allowbreak{}\_\allowbreak{}ActiveKeyId} na jego
        id. Nowe wpisy używają aktywnego klucza, a przy sprawdzaniu próbuje się
        \emph{każdego} wymienionego klucza --- dzięki temu supresja duplikatów nie
        zeruje się.
  \item Stary klucz usuń z konfiguracji dopiero \textbf{po oknie rejestru}
        (domyślnie 365 dni, \texttt{Submission:\allowbreak{}Retention:\allowbreak{}LedgerWindowDays}).
        Wpisy pod nieznanym id są od tej pory nieczynne, aż zbierze je czyszczenie.
\end{enumerate}

Rotacja jest pracą operatora, nie jest zautomatyzowana. Po oknie to samo konto może
znów zgłosić dla tego pracodawcy („jedno na pracodawcę” nie jest gwarancją trwałą);
usunięcie rekordu po paragonie \emph{nie} czyści wpisu w rejestrze. Operator z bazą
\emph{i} kluczem może wyliczyć tagi (T-08, zaakceptowane) --- dlatego klucz trzyma się
\textbf{poza domeną kopii zapasowej bazy}. Rotacji \emph{nie} wykonywano w tej sesji.

\subsection{Retencja i kopie zapasowe}

Parametry czyszczenia (\repopath{docs/architecture/submission-flow.md}): serwis
\texttt{RetentionService} co \texttt{Submission:\allowbreak{}Retention:\allowbreak{}SweepIntervalMinutes}
(5) usuwa w partiach po 500; domyślnie \texttt{RecordMaxAgeMonths} = 24,
\texttt{LedgerWindowDays} = 365; bilety: TTL 15 min, do 3 niewykorzystanych, 10
wydań na godzinę; usuwanie po paragonie: 6 na minutę na IP i 60 globalnie;
zgłoszenie z biletem: 6 na minutę na IP i 60 globalnie.

\textbf{Kopie zapasowe: w repozytorium nie ma żadnej polityki ani procedury.}
Projekt prywatności zapisuje to jako założenie (DESIGN.md §retencja: „no backup
policy exists yet (nothing is deployed)”; usunięcia docierają do kopii
dopiero, gdy kopie się starzeją), a lista kontrolna gotowości do wdrożenia w
\repopath{docs/OWNER-FOLLOWUPS.md} ma otwarte pole „Backup and recovery tested”.
Tego tekstu nie wolno czytać jako procedury odtwarzania.

\subsection{Decyzje zostawione właścicielowi}

Otwarte pozycje z \repopath{docs/OWNER-FOLLOWUPS.md}: topologia jedno- czy
wieloinstancyjna (limity i wydawca sygnałów żyją w pamięci procesu), limity
szybkości na endpointach anonimowych (OP-15), przypięcie authservice po skrócie,
testy na prawdziwym authservice (OP-17), z prawdziwym Claude (OP-18) i przez
prawdziwy TLS (OP-25), przegląd prawnika. Do tego czasu: żadnych prawdziwych
wywiadów.

\section{Rozwiązywanie problemów}

Tabela łączy rzeczy \emph{znane} z repozytorium: błędy z
\repopath{scripts/README.md} (klucz: dosłowny tekst na ekranie), z runbooka
Claude, z dokumentu zgłoszeń CLI oraz te zaobserwowane w tej sesji (oznaczone
\emph{(obserwowane)}). Niczego tu nie wymyślono.

\begin{center}
\scriptsize
\begin{longtable}{@{}>{\raggedright\arraybackslash}p{0.29\linewidth}>{\raggedright\arraybackslash}p{0.32\linewidth}>{\raggedright\arraybackslash}p{0.32\linewidth}@{}}
\toprule
Objaw & Przyczyna & Działanie \\
\midrule
\endhead
\texttt{Container runtime 'docker' could not be found.} / \texttt{Cannot connect to the Docker daemon} / \texttt{failed to connect to the docker API} & brak silnika w \texttt{PATH} albo demon nie działa & zainstaluj Docker lub Podman (\texttt{DOTNET\_\allowbreak{}ASPIRE\_\allowbreak{}CONTAINER\_\allowbreak{}RUNTIME=podman}), uruchom demona; sprawdź \texttt{docker info}. \emph{(obserwowane:} \texttt{setup.\allowbreak{}sh --check} zgłasza \texttt{MISSING container engine} także gdy binarka \texttt{docker} istnieje, a demon nie.) \\
\texttt{scan-\allowbreak{}secrets: neither 'gitleaks' nor a running Docker daemon is available.} & hak pre-commit nie może działać & zainstaluj gitleaks albo uruchom Docker \\
\texttt{initdb: error: directory /var/\allowbreak{}lib/\allowbreak{}postgresql/\allowbreak{}data exists but is not empty} (Postgres wychodzi, AppHost nie jest zdrowy) & wolumen \texttt{exit-\allowbreak{}interview-\allowbreak{}agent-\allowbreak{}pgdata} zainicjowano inną wersją główną Postgresa (AppHost przypina 17) & \texttt{docker volume rm exit-\allowbreak{}interview-\allowbreak{}agent-\allowbreak{}pgdata} (tylko dane deweloperskie) i uruchom ponownie \\
nie da się pobrać \texttt{ghcr.\allowbreak{}io/\allowbreak{}konradcinkusz/\allowbreak{}authservice} & blobs obrazu niedostępne przez proxy (opisane dla piaskownicy w scripts/README) & \texttt{Identity\_\allowbreak{}\_\allowbreak{}Enabled=false dotnet run --project src/\allowbreak{}ExitInterviewAgent.\allowbreak{}AppHost}; punkty chronione dadzą 401 \\
authservice kończy pracę z wzmianką o \texttt{Jwt:Algorithm} & klucz podpisu nie dotarł do authservice, wybrano ścieżkę symetryczną & uruchom \texttt{scripts/setup.*} albo pozwól AppHost wygenerować klucz; sprawdź, że \texttt{GET <authservice>/\allowbreak{}.well-\allowbreak{}known/\allowbreak{}jwks.\allowbreak{}json} zwraca niepustą tablicę \texttt{keys} \\
\texttt{IDX10500:\allowbreak{}Signature validation failed. No security keys were provided} (log interview-service) & brak JWKS: puste \texttt{Jwt:Authority} albo authservice nieosiągalny & uruchom przez AppHost (wstrzykuje autorytet) lub ustaw \texttt{Jwt\_\allowbreak{}\_\allowbreak{}Authority} \\
usługa nie startuje poza Development, brak klucza rejestru & brak \texttt{Ledger\_\_*} & ustaw klucz HMAC (SECRETS.md); poza Development nie ma klucza domyślnego \\
\texttt{Unknown option '--help'.} przy \texttt{interview --help} \emph{(obserwowane)} & podpolecenia nie mają własnej pomocy & użyj \texttt{exit-\allowbreak{}interview --help} (jedna, globalna) \\
\texttt{Usage: exit-\allowbreak{}interview providers [ping ...]} przy \texttt{providers --provider ...} \emph{(obserwowane)} & flagi dostawcy są dozwolone tylko po \texttt{ping} & \texttt{exit-\allowbreak{}interview providers ping --provider ... --model ...} \\
\texttt{The config file must not contain a key or token.} \emph{(obserwowane)} & w pliku konfiguracyjnym jest klucz & usuń go, zostaw \texttt{apiKeyEnv} z nazwą zmiennej \\
\texttt{--save-\allowbreak{}transcript needs --out <dir>} & zapis transkrypcji bez katalogu & dodaj \texttt{--out} \\
\texttt{No server address. Pass --server <url> or set EXIT\_\allowbreak{}INTERVIEW\_\allowbreak{}SERVER\_\allowbreak{}URL.} \emph{(obserwowane)} & nie ma wartości domyślnej, celowo & podaj adres wdrożenia, które masz na myśli \\
\texttt{Plain http:\allowbreak{}/\allowbreak{}/\allowbreak{}is accepted only for a service on this computer} \emph{(obserwowane)} & adres http poza localhost & użyj https \\
\texttt{ALREADY\_\allowbreak{}SUBMITTED} (409) & jedno zgłoszenie na konto i pracodawcę w oknie rejestru; usunięcie nie otwiera ponownie & zaczekaj do końca okna rejestru (domyślnie 365 dni) \\
portal \texttt{/connect}: „The connector is not set up in this environment” & brak \texttt{MCP\_\allowbreak{}RESOURCE\_\allowbreak{}URL} (\texttt{/api/config}: \texttt{mcp.url:null}) & ustaw oba adresy MCP (sekcja~\ref{sec:mcp}) \\
403 \texttt{ORIGIN\_\allowbreak{}NOT\_\allowbreak{}ALLOWED} & wywołanie z przeglądarki & zostaw \texttt{Mcp:\allowbreak{}AllowedOrigins} puste dla Claude \\
\bottomrule
\end{longtable}
\end{center}

\section{Ściągawka}

\subsection{Komendy}

\begin{center}
\scriptsize
\begin{longtable}{@{}>{\raggedright\arraybackslash}p{0.52\linewidth}>{\raggedright\arraybackslash}p{0.2\linewidth}>{\raggedright\arraybackslash}p{0.2\linewidth}@{}}
\toprule
Komenda & Co robi & Status w tej sesji \\
\midrule
\endhead
\texttt{scripts/\allowbreak{}setup.\allowbreak{}sh [--check]} & przygotowanie, raport braków & \texttt{--check} uruchomione; pełne: nie \\
\texttt{dotnet build -warnaserror \&\& dotnet test} & build i testy .NET & 0 ostrzeżeń; 1620 zaliczonych, 17 pominiętych \\
\texttt{scripts/\allowbreak{}check-\allowbreak{}kernel-\allowbreak{}size.\allowbreak{}sh} & limit jądra (800 linii) & 320 linii \\
\texttt{pnpm lint \&\& pnpm typecheck \&\& pnpm test \&\& pnpm build} & portal & uruchomione, 253 testy \\
\texttt{cd tests/\allowbreak{}e2e \&\& pnpm test} (z \texttt{PLAYWRIGHT\_\allowbreak{}CHROMIUM\_\allowbreak{}EXECUTABLE}) & testy przeglądarkowe & 115 zaliczonych (na atrapie) \\
\texttt{dotnet run --project src/\allowbreak{}ExitInterviewAgent.\allowbreak{}AppHost} & cały stos lokalny & \textbf{nieprzetestowane} (brak demona kontenerów) \\
\texttt{Identity\_\allowbreak{}\_\allowbreak{}Enabled=false dotnet run --project ...AppHost} & stos bez authservice & \textbf{nieprzetestowane} \\
\texttt{dotnet user-\allowbreak{}secrets set Mcp:\allowbreak{}ResourceUrl ... --project ...AppHost} & adres MCP & \textbf{nieprzetestowane} \\
\texttt{exit-\allowbreak{}interview --help} / \texttt{--version} & pomoc, wersja protokołu & uruchomione \\
\texttt{exit-\allowbreak{}interview personas} & lista 8 rozmówców & uruchomione \\
\texttt{exit-\allowbreak{}interview demo --persona <id> [--seed n] [--out d]} & wywiad offline na mocku & uruchomione, deterministyczne \\
\texttt{exit-\allowbreak{}interview providers} & konfiguracja bez sieci & uruchomione \\
\texttt{exit-\allowbreak{}interview providers ping --provider p --model m} & jedno żądanie na żywo & uruchomione tylko do niedostępnego Ollama \\
\texttt{exit-\allowbreak{}interview providers forget-\allowbreak{}confirmations} & kasuje zapamiętane zgody & uruchomione \\
\texttt{exit-\allowbreak{}interview interview --provider mock --model mock --tenure 1y\_\allowbreak{}3y --out d} & wywiad offline ze stdin & uruchomione \\
\texttt{exit-\allowbreak{}interview interview --provider anthropic|openai-\allowbreak{}compatible|ollama ...} & prawdziwy wywiad & ujawnienie sprawdzone; sam wywiad \textbf{nieprzetestowany} \\
\texttt{exit-\allowbreak{}interview submit --record f --server u [--save-\allowbreak{}receipt f] [--yes]} & wysłanie rekordu & podgląd i odmowy sprawdzone; sukces \textbf{nieprzetestowany} \\
\texttt{exit-\allowbreak{}interview delete-\allowbreak{}receipt --server u} & usunięcie po paragonie & kontrola kodu sprawdzona; 204 \textbf{nieprzetestowane} \\
\texttt{exit-\allowbreak{}interview tiles --record f [--provider p --model m] [--out d] [--format text|html|json]} & kafelki z rekordu, ADR-0074 & mock i kody 0, 2, 3, 4 sprawdzone; prawdziwy model \textbf{nieprzetestowany} \\
\texttt{scripts/\allowbreak{}run-\allowbreak{}evals.\allowbreak{}sh [validate|gate|report|calibrate]} & harness offline & uruchomione w całości \\
\texttt{dotnet run --project src/\allowbreak{}ExitInterviewAgent.\allowbreak{}Eval -- profiles} & dostępne profile & uruchomione \\
\texttt{... Eval -- run --profile <nazwa> --out report/\allowbreak{}} & raport; prawdziwy model & mock: tak; prawdziwy: \textbf{nie} \\
\texttt{curl localhost:\allowbreak{}5200/\allowbreak{}health} & stan integracji & uruchomione (usługa samodzielnie) \\
\texttt{fly secrets set -a <app> K=v --stage} & sekret we wdrożeniu & \textbf{nic nie wdrożono} \\
\texttt{dotnet test tests/\allowbreak{}ExitInterviewAgent.\allowbreak{}InterviewService.\allowbreak{}Tests --filter "FullyQualifiedName\textasciitilde{}InterviewEndpointTests"} & sesja webowa na mocku: zgoda, rozmowa, wynik, usunięcie (część~6) & uruchomione; 24 zaliczone \\
\texttt{dotnet test tests/\allowbreak{}ExitInterviewAgent.\allowbreak{}InterviewService.\allowbreak{}Tests --filter "FullyQualifiedName\textasciitilde{}Billing"} & kredyty, zakup na atrapie, webhook (część~6) & uruchomione; 52 zaliczone, 3 pominięte (wymagają PostgreSQL) \\
\texttt{scripts/\allowbreak{}render-\allowbreak{}diagrams.\allowbreak{}sh p5-webapp} & rysunek architektury webowej (część~6) & uruchomione \\
\bottomrule
\end{longtable}
\end{center}

\subsection{Cztery rzeczy, o których nie wolno zapomnieć}

\begin{enumerate}[leftmargin=*,itemsep=2pt]
  \item \finding{Nie wchodzi tu żaden prawdziwy wywiad.} Brak polityki
        prywatności, podstawy prawnej, działającego usuwania i przeglądu prawnika
        (README; \texttt{docs/\allowbreak{}legal/\allowbreak{}CONSIDERATIONS.\allowbreak{}md} §4). Używaj wymyślonych
        odpowiedzi i fikcyjnego pracodawcy.
  \item Kod paragonu jest pokazany \emph{raz} i nie da się go odzyskać; nikt nie
        wyszuka zgłoszeń użytkownika.
  \item Klucze API tylko w zmiennych środowiskowych; bilet i kod paragonu tylko
        przez zmienną, ukryty monit lub stdin --- nigdy w argumencie.
  \item Zielony \texttt{mock} i zielone e2e na atrapie nie dowodzą jakości
        prawdziwego modelu ani działania wdrożenia: tego w tej sesji nie
        sprawdzono.
\end{enumerate}
}
